\chapter{广告}
\label{chap:rs-advertising}

用户搜索“防水通勤背包，价格不超过500元”时，平台可能同时找到自然商品和付费推广的商品。一个广告即使与查询相符，也未必应该展示：商品可能无法配送，广告主可能已用完预算，素材可能夸大防水性能，或者这次展示虽然很容易带来一次被记账的购买，却几乎不会增加原本的购买量。广告系统需要回答的不是单一的“用户会不会点击”，而是怎样在用户需求、广告主价值与有限展示机会之间完成匹配，并取得能够指导下一轮投放的证据。

前面的章节已经建立候选、响应、展示和受约束决策的接口。第~\ref{chap:rs-search}~章补充查询相关性，本章则加入资金和市场：响应估计提供价值依据，广告主或其控制器提交出价，平台决定分配与支付，实际支出改变下一次请求的资格。创意决定用户实际看到什么，效果测量决定一次成交能够支持怎样的投放判断。这些职责相互依赖，却不能互相替代。预测准确不等于支付规则正确，用完预算不等于取得最大价值，归因收入增加也不必然意味着广告增加了收入。

\section{广告任务与价值}
\label{sec:rs-advertising-task}

\term{赞助搜索}（sponsored search）是在查询结果中分配付费展示机会，查询为匹配提供显式需求；\term{展示广告}（display advertising）使用页面、应用或广告位情境以及允许使用的受众信息匹配；\term{信息流广告}（in-feed advertising）还要与自然内容共同组成连续浏览列表。这些形态可以共享响应模型，却有不同的位置、可见性和交互条件。用户希望看到相关且不误导的内容，广告主关注成交、客户或品牌目标，平台同时面对收入、体验和供给关系。任何单一点击率都不足以代表这三方目标。

记时点$t$的一次机会为$(\bm{x}_t,j,C_t,H_t)$：$\bm{x}_t$包含请求前可用的信息，$j$是待分配的位置，$C_t$是通过当前资格核验的广告标识集合，$H_t$包含预算余额、频次和投放状态。一个广告标识还要关联商品、广告主、广告组和素材版本。同一商品的两个标题版本不能在日志中丢失版本区别；两个渠道中的同一广告主也不能因为标识不同就重复使用共同预算。这里的机会只是可供竞争的请求，不是已经发生的曝光。

为固定响应分母，设$E=1$表示广告实际曝光，$C=1$表示用户点击该广告，$Z=1$表示点击后在约定窗口$W$内完成一次目标转化，未点击时规定$Z=0$。在给定广告、位置及请求信息后，定义
\begin{equation}
  \begin{aligned}
    q&=\mathbb{P}(C=1\mid E=1,\bm{x},i,j),\\
    h&=\mathbb{P}(Z=1\mid C=1,E=1,\bm{x},i,j),\\
    z&=\mathbb{P}(C=1,Z=1\mid E=1,\bm{x},i,j)=qh.
  \end{aligned}
  \label{eq:rs-advertising-chain}
\end{equation}
$q$是点击率（CTR），$h$是\term{点击后转化率}（post-click CVR），$z$是\term{经点击的展示到转化率}（CTCVR）。乘积来自条件概率的链式法则，不要求点击和转化独立。恰恰相反，$h$已经以点击为条件。一次点击有多次订单时，应另建转化次数或金额目标，不能让一个二元标签同时表示“至少一次转化”和“订单总数”。

\term{浏览后转化}不能默默放进$z$。例如，定义$Z^{\mathrm{view}}=1$为曝光后窗口内发生目标转化、但没有点击本广告，并按既定规则归于这次浏览。若两条链路按“是否点击本广告”互斥，则浏览链路的概率为$(1-q)h^{\mathrm{view}}$，其中$h^{\mathrm{view}}=\mathbb{P}(Z^{\mathrm{view}}=1\mid C=0,E=1,\bm{x},i,j)$。若报表允许同一转化同时进入点击与浏览口径，两项不能直接相加。跨设备未匹配到的点击也不能自动证明用户没有点击。

\begin{example}[固定窗口的广告事件账目]
\label{ex:rs-advertising-events}
以下均为教学假设。对1000次成熟曝光观察到100次点击、20次经点击转化，故经验CTR为$100/1000=0.1$，点击后CVR为$20/100=0.2$，CTCVR为$20/1000=0.02$。若其余900次曝光中有9次符合浏览归因规则的转化，浏览条件转化率为$9/900=0.01$，对应全部曝光的比例为0.009，而非0.01。只有在转化互斥、窗口一致且记录完整时，两条链路合计才是0.029。该账目描述观察及核算结果，尚未回答不展示广告会发生多少转化。
\end{example}

事件概率还不是金额。设$M$为转化发生时的约定价值，例如订单收入或扣除商品履约成本后的贡献毛利。仅考虑点击链路，一次曝光的预测价值为
\begin{equation}
  \widehat v^{\mathrm{imp}}
  =\widehat q\,\widehat h\,
    \widehat{\mathbb{E}}[M\mid C=1,Z=1,E=1,\bm{x},i,j].
  \label{eq:rs-advertising-value}
\end{equation}
这一步使用事件条件下的金额，不是把全体订单的平均客单价无条件乘给每次机会。延续例~\ref{ex:rs-advertising-events}，假设每次转化的贡献毛利均为100元，则每次曝光的预测毛利为2元，每次点击的预测毛利为20元。如果分子改成400元的销售收入，两种价值分别变为8元和80元；这不是模型性能变化，而是价值定义变化。

计费方式决定何时产生支付义务。\term{每千次展示成本}（CPM）的报价单位是元/千次展示，比较单次曝光时须除以1000；\term{每次点击成本}（CPC）在有效点击时计费；\term{每次行动成本}（CPA）按约定、可核验且可能迟到的行动计费。优化目标为转化并不意味着按转化收费。目标CPA或优化CPC投放仍可在每次点击扣费，只是控制器按转化价值调整点击出价。

\begin{table}[htbp]
  \centering
  \begin{tabularx}{\linewidth}{@{}lX@{}}
    \toprule
    量 & 含义与账目边界\\
    \midrule
    $v_{t,i}$、$\widehat v_{t,i}$ & 真实私人价值及其估计；须注明按曝光、点击还是行动计量。\\
    $b_{t,i}$ & 提交给机制的出价，不必等于真实价值或价值估计。\\
    $\tau_{t,i}$ & 机制确定的每次计费事件支付额；未必等于出价。\\
    $s_{t,i}$ & 实际计费事件对应的支出；CPC下是有效点击指示量乘$\tau_{t,i}$。\\
    $S_t$、$B_t$ & 截至$t$的累计实际支出与预算余额；无调整时$B_t=B-S_t$。\\
    $\gamma_{k,i}$ & 第$k$次已发生转化分给触点$i$的信用，与付款记录不同。\\
    \bottomrule
  \end{tabularx}
  \caption{广告价值、交易与效果核算使用不同的量。预测价值不能当作实际扣费，转化信用也不能作为未观察到的增量。}
  \label{tab:rs-advertising-ledger}
\end{table}

图~\ref{fig:rs-advertising-loop}把这些量放回执行过程。平台拥有分配和支付规则，广告主或代理控制器选择出价与参与，周期控制器管理资金使用。一个出价乘子只改变报价，参与概率只改变是否竞价，广告组预算限制该组累计支出，渠道预算则限制交给渠道的资金。它们作用范围不同，不能用降低平均出价代替对每次超预算交易的拒绝。

\begin{figure*}[htbp]
  \centering
  % Adapted from the academic-drawing TikZ template to the book's palette and fonts.
\begin{tikzpicture}[
  font=\figurefont\small,
  box/.style={rectangle,draw=BookInk,line width=0.85pt,
    minimum height=12mm,text width=25mm,align=center,inner sep=2mm},
  flow/.style={-{Latex[length=2.2mm,width=1.5mm]},draw=BookInk,
    line width=1.1pt,rounded corners=1.2mm},
  feedback/.style={flow,dashed,draw=BookAmber,line width=0.85pt}
]
  \node[box,fill=FigureInput!12] (eligible) at (0,0) {请求与资格\\合格候选};
  \node[box,fill=FigureModel!12] (value) at (4.1,0) {响应与价值\\$\widehat q,\widehat h,\widehat v$};
  \node[box,fill=FigureModel!12] (auction) at (8.2,0) {出价、分配\\与事件支付};
  \node[box,fill=FigureOutput!12] (execute) at (12.3,0) {曝光与结算\\实际支出$s_t$};
  \node[box,fill=BookPaper] (budget) at (8.2,-2.4) {周期控制\\余额与预留};
  \node[box,fill=BookPaper] (learn) at (4.1,-2.4) {标签与测量\\成熟、信用、增量};
  \draw[flow] (eligible.east) -- (value.west);
  \draw[flow] (value.east) -- (auction.west);
  \draw[flow] (auction.east) -- (execute.west);
  \draw[feedback] (execute.south) |- (budget.east);
  \draw[feedback] (budget.north) -- (auction.south);
  \draw[feedback] (learn.north) -- (value.south);
  \draw[feedback] (execute.south) -- (12.3,-3.9)
    -| (learn.south);
\end{tikzpicture}

  \caption{广告执行与反馈的职责示意，并非所有平台的固定部署顺序。实线传递当前候选、预测与执行结果；虚线传递跨请求状态和学习反馈。分配与支付共同遵守展示可行域，不能在拍卖结束后随意更换获胜者却沿用原支付。}
  \label{fig:rs-advertising-loop}
\end{figure*}

目标成本与回报也需要明确分子。对同一周期，记转化数为$N_{\mathrm{conv}}$，销售收入为$R$，扣除商品成本但尚未扣广告费的贡献毛利为$G$，广告实际支出为$S$，则
\begin{equation}
  \operatorname{CPA}=\frac{S}{N_{\mathrm{conv}}},\qquad
  \operatorname{ROAS}=\frac{R}{S},\qquad
  \operatorname{ROI}_{\mathrm{profit}}=\frac{G-S}{S}.
  \label{eq:rs-advertising-returns}
\end{equation}
若每次转化收入400元、贡献毛利100元，20次转化花费1000元，则CPA为50元，\term{广告支出回报}（ROAS）为8，\term{利润ROI}为1。三者不矛盾。文献中的ROI、ROS有时实际指约定价值与支出的比值，本章比较时以公式为准。$S=0$或$N_{\mathrm{conv}}=0$时不应强行报一个有限比值；优化时可用$S\leq cN_{\mathrm{conv}}$或$\rho S\leq V$表达目标CPA和目标价值回报，但必须说明$V$是预测、归因还是增量价值。

预算是跨请求状态。若先前已支出60元，总预算为100元，本次有效点击的支付是3元，则新余额是37元，与本次出价是4元还是5元无关。点击和行动计费可能在展示后回流，执行时还要为在途义务保留资金或明确允许的账期超支政策。硬预算核验使用可用余额扣除预留额，最终结算则用实际支付冲销预留。并发请求中的“每个请求都看见足够余额”也不构成总预算保证，需要原子预留或等价的额度协调。

频次具有相似的状态性，但消耗单位是曝光而不是货币。第~\ref{sec:rs-presentation-time}~节已经解释商品、素材和品牌频次，本章把频次上限接入资格判定：一个广告主即使还有资金，当前用户也可能已达到窗口内的曝光上限。学习到的疲劳响应可以降低价值估计，却不能代替硬频控。响应标签、资格快照、素材版本与预算状态均须按请求时点保存，后文才能区分标签未成熟、资格变化与策略造成的选择。

\section{广告匹配与响应预估}
\label{sec:rs-advertising-matching}

广告投放需要两类输入：有哪些广告能够参与，以及在指定展示条件下它们可能产生什么响应。前者检验资格和覆盖，后者估计概率与价值。只有先固定广告、素材、位置及标签窗口，两项任务的结果才有可比性；把不合格广告的高分保留下来并不能增加有效覆盖。

\subsection{广告资格与匹配}
\label{sec:rs-advertising-eligibility}

赞助搜索可以使用关键词或语义条件，展示和信息流可以使用允许的地域、时间、页面与受众定向。匹配成功之后仍要检查商品库存、配送区域、素材审核、广告主状态和当前额度。例如，“防水通勤背包”不能因为文本相似就匹配到仅有防泼水涂层的商品；落地页价格超过500元时也不能继续借用旧标题声称满足价格上限。候选接口沿用第~\ref{chap:rs-retrieval}~章，查询解释沿用第~\ref{chap:rs-search}~章。

资格过滤可以前移到索引分区，也可以在取回候选后精确核验。前移节约后续计算，却可能因缓存过旧漏掉刚恢复的供给；后置更接近当前事实，却可能让大量不可投放候选占满截取名额。工程上应同时记录过滤前覆盖、合格后覆盖和请求代价，不能只看原始召回数量。预算与频控由对应状态服务提供，不由检索分数猜测。

\subsubsection{个性化广告检索框架}
\label{sec:rs-advertising-retrieval}

只靠广告主提前选择的关键词，相关广告可能因为没有购买恰当词项而无法进入候选。\term{个性化广告检索框架}把查询、近期点击、长期点击和画像作为信号，经由中间检索键到达广告。它的“信号—键—广告”三级网络不是一个逐层训练的图神经网络：信号到键是重写边，键到广告是选广告边，模型学习哪些边及其组合值得保留\scite{rs-a01}

键可以是查询、商品、店铺或品牌。选择键要兼顾相关性和倒排长度：过于宽泛的年龄类别可能连到大量无关广告，过细的孤立商品又缺少覆盖。原文以点击计数、结合点击与曝光占比的修正信息值，以及同会话动作的余弦相关性建立初始边。点击计数提供联系证据，却会偏向曝光多的广告；加入曝光统计有助于区分“经常出现”和“出现后容易点击”，会话证据则补充共同需求下的可达路径。初始图仍受历史覆盖限制，后续学习无法从完全不存在的边中恢复一个候选。

一个训练样本由请求信号集合、实际展示广告和点击标签$c_n$组成。沿初始网络找到所有激活的重写边与选广告边，把边ID、点击数、曝光数和历史CTR组织成特征$\bm{\phi}_n$。例如，逻辑回归的点击目标为
\begin{equation}
  \begin{gathered}
    \min_{\bm{w},a}
    \sum_{n=1}^{N}
      \ell\!\left(c_n,\sigma(a+\langle\bm{w},\bm{\phi}_n\rangle)\right)
      +\lambda\lVert\bm{w}\rVert_2^2,\\
    \ell(y,p)=-y\log p-(1-y)\log(1-p).
  \end{gathered}
  \label{eq:rs-advertising-retrieval-loss}
\end{equation}
统计特征只能由样本之前的数据计算。小批量梯度更新边权和统计特征系数，独立时段选择正则、边保留阈值和候选数。原文也比较GBDT与MLP；LR同时使用稀疏ID及连续统计，后两者采用连续特征。树的分裂或MLP的非线性交互改变评分能力，线上仍须承担相应计算，不能因离线更复杂就假设检索代价不变。

若正点击样本按广告价格$a_{\mathrm{price}}>0$加权、负例保持权重1，同一条件组中的最优输出变为
\[
  p_{\mathrm{weighted}}
  =\frac{a_{\mathrm{price}}q}{1-q+a_{\mathrm{price}}q}.
\]
这是对价格加权损失的响应，不是原始CTR。教学假设$q=0.1$、价格权重为2，则输出为$2/11\approx0.182$，而真实点击比例仍是0.1。其赔率为$a_{\mathrm{price}}q/(1-q)$；低CTR下与价格乘CTR近似成比例，但不能直接把它当作校准概率、广告主真实价值或最终支付。

训练后把保留的重写边写入重写倒排索引，把保留的选广告边写入选广告倒排索引，并存储相应权重与统计。线上提取请求信号，依次访问两套索引，按广告标识合并路径，再截取候选。原文的线上LR实现对激活边权求和形成检索排序分数；同一广告经多条路径到达，不意味着应多发一次曝光。一个边被不同路径共享时，须遵循既定特征计数规则，不能因遍历实现重复累加本应只激活一次的ID。

沿背包例子，作如下教学设定：查询“通勤背包”到键“电脑背包”的重写边权为0.4，该键到广告甲的选广告边权为0.5；最近点击的轻量背包到品牌键的权重为0.2，该品牌键到甲的权重为0.3、到乙的权重为0.6。四条不同激活边使甲得分$0.4+0.5+0.2+0.3=1.4$，乙由后一条路径得分$0.2+0.6=0.8$。若另一独立路径使丙得分1.0，去重后的顺序是甲、丙、乙。假设甲已缺货、丙可配送且素材有效，则截取前两项后过滤仅剩丙；若目标是取得两个合格候选，需要扩大访问或补取乙，而不是把缺货甲交给评分器。

设本次访问的重写边与选广告边总数为$L$，到达的不同广告数为$M$，用哈希合并及大小为$K$的堆截取时，在线代价约为$O(L+M\log K)$，存储随保留边数增长。这个复杂度尚未包含资格查询和统计特征读取。路径裁剪减少$L$却可能降低相关广告覆盖，热门键带来长倒排，陈旧统计可能偏向过时需求。原文配套的优化CPC价格来自CVR、商品价格和营销比例，第~\ref{sec:rs-advertising-single}~节再说明其计价口径；它不是三级检索网络自动推导出的拍卖价格。

\subsubsection{OKG：广告主关键词生成}
\label{sec:rs-advertising-okg}

广告主维护词表时面对的是另一项任务：商品已有新的用途、价格或卖点，但当前词表可能没有覆盖。\term{OKG}把语言模型与搜索、广告API、向量记忆和计算工具组合，持续产生广告主可采用的关键词集合。搜索补充产品信息，广告API取得已投放词的表现，记忆保存词项、产品事实与反馈，计算工具决定下一轮应扩展哪些类别\scite{rs-a02}

记第$t$轮可核实信息为$F_t$，已有词集合为$K_t$，表现统计为$\bm{m}_t$，一次维护输出$K_{t+1}=g(F_t,K_t,\bm{m}_t)$。这里的“学习”是更新可检索信息、提示与候选状态，并不表示每次API调用后都用梯度修改语言模型参数。原文将固定总量$n$的新词分为扩展新类别和细化已有成功类别两部分。若上一轮两部分的目标KPI分别为$m_{\mathrm{wide}}$和$m_{\mathrm{deep}}$，计算配额为
\begin{equation}
  n_{\mathrm{wide}}
  =\left\lfloor
    n\frac{m_{\mathrm{wide}}}{m_{\mathrm{wide}}+m_{\mathrm{deep}}}
   \right\rfloor,\qquad
  n_{\mathrm{deep}}=n-n_{\mathrm{wide}}.
  \label{eq:rs-advertising-okg-quota}
\end{equation}
比例分配只适用于可加且方向一致的KPI；不能把希望降低的CPC未经变换直接当成越大越好的份额。分母为零时需要预设探索配额，计数窗口、流量规模和归类方式也必须一致。固定新词数量控制维护规模，不是预算不超支的数学保证。

教学假设某背包已核实有15英寸电脑隔层、20升容量和防泼水面料，售价399元。上轮新类别词带来30次点击、细化类别带来70次点击，本轮生成10个词，则新类别3个、细化类别7个。检索产品事实和既有词后，新类别可探索“短途出差背包”“轻量日用背包”等使用需求，细化通勤类别可生成“15英寸电脑通勤包”“399元通勤背包”等。每个候选都要标明依据；“专业登山防水包”“全天候暴雨防水”没有事实支持，应删除，而不是因语言流畅就进入投放。

执行时先取得产品与历史反馈快照，再计算配额，调用模型生成、去重和核验词项，把通过的词交给广告主采用并记录版本。到下一轮，广告API提供的仍是已采用词的表现，而非所有被生成词的潜在效果。无效词、重复词和因产品变化失效的词需要撤回；这一维护循环通常不在用户请求的同步检索路径上。实际请求仍由广告匹配接口执行，不能为了等待网页搜索或长文本生成而阻塞每次曝光机会。

原文的离线评价把生成词与已有离线词做语义近邻匹配，再读取匹配词的KPI代理。因此，“相近旧词曾经有较多点击”与“新词真实投放会增加点击或成交”是两项证据。相似度门槛可能排除生僻但有用的新词，旧词流量与出价也可能不同，代理KPI不能提供因果增量。真实新增覆盖与效果应按第~\ref{sec:rs-advertising-measurement}~节的协议验证。

若每轮调用$J$次工具、生成$L_{\mathrm{text}}$个词元、记忆有$M$条记录，成本由工具等待、检索与模型解码共同决定；全扫描记忆是$O(Md)$，其中$d$为嵌入维度，近似索引则有自己的召回误差。这里没有仅由关键词数$n$决定的固定低时延保证。错误网页、过期价格、广告API缺失反馈和反复强化高流量旧词都会影响词表；事实核验与保留探索份额解决不同问题，不能互相代替。

\subsection{广告响应预估}
\label{sec:rs-advertising-response}

合格候选只保证可以投放，不保证值得投放。响应模型把广告、素材、用户允许使用的历史、位置与请求情境映射为式~\eqref{eq:rs-advertising-chain}中的概率。模型输入应保留实际素材版本，目标应保留点击或转化事件及窗口；不同版本共用商品表示，不等于它们拥有相同点击率。广告价格和预算可以参与决策，但不能成为从未来账目回填到训练样本的特征。

\subsubsection{观测数据下的响应学习}
\label{sec:rs-advertising-observed-response}

特征交互、点击选择与负例采样沿用第~\ref{sec:rs-observed-learning}~节。DIN、DIEN与SIM分别复用第~\ref{sec:rs-din}~节、第~\ref{sec:rs-dien}~节及第~\ref{sec:rs-sim}~节的候选相关兴趣、兴趣演化和长历史访问：这里把待评分广告及其素材作为目标对象，输出指定位置的CTR。SIM取回的是用于解释当前广告的历史行为，不是本次广告候选。MetaEmb沿用第~\ref{sec:rs-metaemb}~节初始化新广告嵌入；新素材的少量点击不能当成新商品的全部转化证据，初始化也不能缩短转化窗口。

广告中容易同时出现三种困难。点击过的样本不代表全部曝光，转化远少于点击，而新近点击的转化标签还未成熟。ESMM和ESCM$^2$主要组织前两种条件下的学习；延迟方法处理第三种条件。把两类方法组合前，必须先说明反事实风险使用的是成熟转化标签还是经过延迟修正的风险，不能简单相加损失后就宣称同时消除了全部偏差。

\paragraph{ESMM：全空间联合响应估计}
\label{sec:rs-advertising-esmm}

仅用点击样本训练CVR，既缺少转化监督，也无法利用大量未点击曝光的表示信息。\term{全空间多任务模型}（Entire Space Multi-task Model，\term{ESMM}）保留全部曝光，为CTR和CVR设置两个分支，共享输入嵌入，再把两分支输出相乘得到CTCVR\scite{rs-a03}

训练集是$D=\{(\bm{x}_n,c_n,z_n)\}_{n=1}^{N}$，每项为一次曝光；$z_n=1$必须同时满足点击及窗口内转化，未点击时为0。设$\widehat q_n=f_{\mathrm{ctr}}(\bm{x}_n)$、$\widehat h_n=f_{\mathrm{cvr}}(\bm{x}_n)$，两者均通过sigmoid限制在$[0,1]$内。利用式~\eqref{eq:rs-advertising-retrieval-loss}的二元交叉熵，联合目标为
\begin{equation}
  L_{\mathrm{ESMM}}
  =\frac{1}{N}\sum_{n=1}^{N}
    \left[
      \ell(c_n,\widehat q_n)
      +\ell(z_n,\widehat q_n\widehat h_n)
    \right].
  \label{eq:rs-advertising-esmm}
\end{equation}
两个损失的分母都是曝光数$N$。第二项的梯度经乘积流入CVR分支，不另把未点击曝光标成“点击后未转化”训练一个CVR交叉熵。这一区别决定了标签含义：未点击的$z_n=0$是CTCVR负例，不是已观察到的点击后CVR负例。

小批量训练时更新共享嵌入和两个分支；共享层从更密集的点击任务获得信号，CVR分支从转化链路获得监督。线上对合格广告计算$\widehat q_n,\widehat h_n$及其乘积，可以按CPC价值使用$\widehat h_n$，也可以按曝光价值使用$\widehat q_n\widehat h_n$。若CTR及CTCVR各自独立训练后再相除，小CTR会放大误差，且商可能超过1；乘积参数化直接避免了后一问题。

在例~\ref{ex:rs-advertising-events}的同质教学数据上，$\widehat q=0.1$、$\widehat h=0.2$给出$\widehat z=0.02$。CTR项平均损失为$-0.1\log0.1-0.9\log0.9\approx0.3251$，CTCVR项为$-0.02\log0.02-0.98\log0.98\approx0.0980$。100次点击中的20次转化支持CVR为0.2，而不是用1000次曝光作CVR分母得到0.02。

设共享表示及两个分支每样本前向代价分别为$C_{\mathrm{emb}}$、$C_q$、$C_h$，一轮训练约为$O(N(C_{\mathrm{emb}}+C_q+C_h))$，额外乘积仅为常数代价。共享参数不保证两任务完全相容，较差的CTR会把误差传给CVR；有限样本、模型失配和小点击倾向也会限制条件概率估计。全空间是训练样本范围，不是“任意日志下CVR无偏”的证明。式~\eqref{eq:rs-advertising-esmm}也没有处理尚未来到的转化，更不包含浏览后转化。

% 标题、导语与完整边注文献保留在同一页。
\Needspace{90mm}
\paragraph{ESCM$^2$：反事实多任务响应估计}
\label{sec:rs-advertising-escm2}

ESMM从乘积任务间接学习CVR，但业务还希望知道：对同类曝光机会，假如发生点击，转化风险如何？\term{ESCM$^2$}保留CTR与CTCVR任务，并加入直接面向CVR的反事实风险正则\scite{rs-a04}这里要区分观测条件概率$h(\bm{x})$与“令点击发生”的潜在转化概率。只有在一致性、给定处理前特征后点击选择可忽略、且点击倾向非零等条件下，二者才可对应；模型名称本身不提供这些条件。

对点击曝光，记成熟转化标签为$y_n$，风险为$\delta_n=\ell(y_n,\widehat h_n)$；未点击曝光没有这个事实标签。设$e_n=\mathbb{P}(C=1\mid\bm{x}_n,E=1)$为点击倾向，由CTR分支估计$\widehat e_n$。逆倾向版本对曝光空间的CVR风险使用
\begin{equation}
  \widehat R_{\mathrm{IPS}}
  =\frac{1}{N}\sum_{n=1}^{N}
    \frac{c_n}{\widehat e_n}\delta_n,\qquad
  L=L_{\mathrm{ctr}}+\lambda_{\mathrm{g}}L_{\mathrm{ctcvr}}
       +\lambda_{\mathrm{c}}\widehat R_{\mathrm{IPS}}.
  \label{eq:rs-advertising-escm2-ips}
\end{equation}
点击倾向低的已点击样本代表更多同类曝光，权重因此较大。这里估计的是曝光分布下的潜在CVR风险，不是简单重算已点击子集的平均损失。联合损失通过共享表示与分支参数求梯度，$\lambda_{\mathrm{g}}$、$\lambda_{\mathrm{c}}$在独立时段验证，不能由测试集效果反复选择。

双重稳健版本另训练误差补全分支$\widehat\delta_n=m(\bm{x}_n)$，估计CVR预测的损失，而不是再把某个转化概率直接当作损失。风险及其辅助训练目标可写成
\begin{equation}
  \begin{aligned}
    \widehat R_{\mathrm{DR}}
      &=\frac{1}{N}\sum_n
        \left[\widehat\delta_n+
          \frac{c_n}{\widehat e_n}
            (\delta_n-\widehat\delta_n)\right],\\
    L_{\mathrm{imp}}
      &=\frac{1}{N}\sum_n
        \frac{c_n}{\widehat e_n}
        (\delta_n-\widehat\delta_n)^2.
  \end{aligned}
  \label{eq:rs-advertising-escm2-dr}
\end{equation}
补全分支从点击样本的损失拟合全部曝光的风险，加权残差修正补全误差。原文将风险项与补全训练项纳入CVR正则；实现时应明确各次更新的参数和停止梯度位置，使补全目标随当前CVR模型更新，又不让辅助模型通过任意改变训练目标逃避拟合。双重稳健性描述在识别条件下，倾向或条件损失模型之一正确时的风险性质，不表示两个模型都错时仍正确，也不消除有限样本方差。

作两组教学假设：甲、乙各100次曝光，点击倾向分别为0.1、0.5，实际点击数恰为10、50；两组已点击样本的平均CVR损失分别为0.2、0.6。直接平均点击损失为$(10\times0.2+50\times0.6)/60=0.5333$。按曝光分布修正后是$(10\times10\times0.2+50\times2\times0.6)/200=0.4$，恰为两组风险的等权平均。若误差补全分别是0.25、0.55，DR为$0.4+(10\times10\times(-0.05)+50\times2\times0.05)/200=0.4$。这是倾向恰好正确的算例，不是对真实数据的保证。

当倾向降至0.001，一个点击权重会达到1000。裁剪倾向、归一化权重或增大训练批量有助于稳定，但前两项改变估计量，须重新报告偏差与有效样本量。若某类广告从不被点击，就没有支撑其反事实CVR的同类事实记录。IPS比ESMM增加常数级权重计算，DR还增加一个误差分支；线上通常只需CTR、CVR及乘积输出，辅助风险模型不必进入每次竞价。最终仍用成熟、未参与训练的点击与曝光数据分别检查CVR和CTCVR。

\paragraph{ES-DFM：经过时间采样}
\label{sec:rs-advertising-esdfm}

转化选择风险明确之后，还要解决“训练时标签是否已经可知”。\term{经过时间采样延迟反馈模型}（Elapsed-Time Sampling Delayed Feedback Model，\term{ES-DFM}）为一次点击选择等待时间，使较新的记录可以较早进入训练，再通过迟到正例回流与权重调整修正训练分布\scite{rs-a06}

以下各延迟方法均以点击为样本单位。记点击时刻为$t_{\mathrm{click}}$，转化延迟为$D$，初次等待为$e$，归因窗口为$W>e$。目标$Y=1$表示$D\leq W$；超过$W$的转化对这个目标为0，不表示此人今后永不购买。每项日志保存请求时特征、点击时间、已知转化时间、当前采样时刻、可见标签$\widetilde y$及回流类型。等待时间可以从$P(e\mid\bm{x})$抽取，也可采用固定窗口，选取依据是历史延迟与新鲜度验证，不能看到该样本未来转化后再为它挑选有利等待时间。

为展示权重来源，暂在固定$\bm{x}$下记真实窗口CVR为$h$、迟到正例概率为$d=\mathbb{P}(Y=1,D>e\mid\bm{x})$，暂时负例中真实负例的比例为$r=(1-h)/(1-h+d)$。初次训练包含比例$h-d$的正例和$1-h+d$的暂时负例；迟到后另补$d$份正例，总记录量变成$1+d$。于是回流后训练记录中的正例概率为$h/(1+d)$，负例概率为$(1-h+d)/(1+d)$。对应权重为
\begin{equation}
  \begin{gathered}
    w_+=1+d,\qquad w_-=(1+d)r,\\
    L_{\mathrm{ESDFM}}
    =\mathbb{E}_{\mathrm{stream}}
      [\widetilde Yw_+\ell_+
        +(1-\widetilde Y)w_-\ell_-],
  \end{gathered}
  \label{eq:rs-advertising-esdfm}
\end{equation}
其中$\ell_+=-\log\widehat h$、$\ell_-=-\log(1-\widehat h)$。这里的$d$是全部点击中迟到正例的概率，不是暂时负例内部的假负例比例。

辅助模型在已成熟日志上学习$d$与$r$：前者把迟到转化标为1、其他点击标为0；后者排除初次窗口内正例，在余下点击中将成熟负例标为1、迟到正例标为0。可用共享网络与掩码训练两个分类头，再把预测作为主模型的权重。流式阶段按到达顺序更新主CVR，权重不应利用尚未发生的转化标签。在线推断仅输出$\widehat h$，不必等待当前用户真正转化。

\begin{figure*}[htbp]
  \centering
  % Timeline adapted from the academic-drawing TikZ template.
\begin{tikzpicture}[
  font=\figurefont\small,
  state/.style={rectangle,draw=BookInk,line width=0.85pt,
    text width=25mm,minimum height=13mm,align=center,inner sep=2mm},
  flow/.style={-{Latex[length=2.2mm,width=1.5mm]},
    draw=BookInk,line width=1.1pt},
  alternate/.style={flow,dashed}
]
  \node[state,fill=FigureInput!12] (click) at (0,0)
    {点击发生\\$t=0$};
  \node[state,fill=FigureModel!12] (wait) at (4.1,0)
    {尚未转化\\$e=10$分钟};
  \node[state,fill=FigureOutput!12] (positive) at (8.2,1.2)
    {发生转化\\$D=2$小时};
  \node[state,fill=BookPaper] (negative) at (8.2,-1.5)
    {窗口内无转化\\$W=24$小时};
  \node[align=left,anchor=west,text width=28mm] at (10.2,1.2)
    {正标签回流\\FN对应DP};
  \node[align=left,anchor=west,text width=28mm] at (10.2,-1.5)
    {成熟负标签\\可作为DN回流};
  \draw[flow] (click.east) -- (wait.west);
  \draw[flow] (wait.east) -- ++(0.4,0) |- (positive.west);
  \draw[alternate] (wait.south) |- (negative.west);
  \node[align=center,text width=33mm] at (4.1,1.5)
    {首次可见标签为0\\此刻分不清FN与RN};
  \node[align=left,anchor=west,text width=132mm] at (-1.4,-3)
    {实线与虚线表示两种互斥后续；没有转化的记录须等待窗口结束，不能提前认定成熟。};
\end{tikzpicture}

  \caption{同一点击的标签可见性与回流记录。教学窗口为$e=10$分钟、$W=24$小时；图中2小时转化的是迟到正例，另一条路径到窗口结束仍未转化，才对该窗口成为成熟负例。图示时间点不按比例绘制。}
  \label{fig:rs-advertising-delay}
\end{figure*}

图~\ref{fig:rs-advertising-delay}中，一次点击10分钟时未转化，以$\widetilde y=0$训练；2小时转化，再以$\widetilde y=1$训练。教学假设一组100次点击含10次窗口内即时转化、10次迟到转化和80次窗口内未转化，则$h=0.2$、$d=0.1$、$r=8/9$。初次100条中有90条负记录，后补10条正记录，$w_+=1.1$、$w_-=44/45\approx0.9778$。按110条记录平均，正损失系数为$20\times1.1/110=0.2$，负损失系数为$90\times(44/45)/110=0.8$。对那一个迟到点击，两次贡献为$(44/45)\ell_-+1.1\ell_+$；它并不逐样本抵消成单个正例，修正成立在给定分布的期望中。

这种推导还涉及特征分布：仅重复迟到正例会改变不同$\bm{x}$出现的次数，原文从条件权重到联合风险使用$P_{\mathrm{stream}}(\bm{x})\approx P(\bm{x})$近似。迟到比例较大或不同广告相差悬殊时，不能略去该条件。辅助模型滞后、追踪漏报、窗口改动也会产生新误差。若原点击数为$N$、迟到比例为$\bar d$，主训练约处理$N(1+\bar d)$条记录，另承担两个辅助头；等待队列及特征快照存储随点击到达率与保留窗口增长。这是以数据和计算换取更及时学习，不是免费获得成熟标签。

\paragraph{DEFUSE：细粒度标签修正}
\label{sec:rs-advertising-defuse}

ES-DFM降低暂时负例的负损失权重，但暂时负例中确有未来正例。\term{DEFUSE}进一步区分即时正例（IP）、假负例（FN）、真实负例（RN）和迟到正例（DP），把FN应当贡献的正损失纳入训练\scite{rs-a07}FN是迟到点击的首次记录，DP是同一点击转化后的回流记录，不是两次独立转化。

沿用$h,d,r$，令$\rho=1-r=d/(1-h+d)$为暂时负例属于FN的概率。IP与DP在流式记录中可区分，FN和RN在初次窗口结束时不可区分。DEFUSE用成熟数据训练$\widehat d$与$\widehat\rho$，后者可由$1-\widehat r$获得，也可用迟到概率与当前CVR构成的比例估计。后一方式包含分母与两个模型的误差，不能默认更稳定。

以原文在ES-DFM回流机制上的设置为例，IP与RN权重为$1+d$，DP权重为1，FN权重为$d$。三种可见记录的风险为
\begin{equation}
  L_{\mathrm{record}}=
  \begin{cases}
    (1+\widehat d)\ell_+, & \text{即时正记录},\\
    \widehat\rho\,\widehat d\,\ell_+
      +(1-\widehat\rho)(1+\widehat d)\ell_-,
      & \text{初次负记录},\\
    \ell_+, & \text{迟到正记录}.
  \end{cases}
  \label{eq:rs-advertising-defuse}
\end{equation}
两步优化先取得辅助概率，再用修正风险更新CVR；成熟日志持续提供辅助监督。初次负记录不再只有负损失，但也没有被武断改成正例。一个可见0标签只告诉系统此刻尚未观察到转化，不能告诉系统它必然属于哪类。

在前述100次点击教学组中，$\rho=1/9$、$d=0.1$。每条初次负记录贡献$(1/90)\ell_++(44/45)\ell_-$。若主模型当前输出0.2，直接负例损失为$-\log0.8\approx0.2231$，修正贡献为$(1/90)(-\log0.2)+(44/45)(-\log0.8)\approx0.2361$。修正并不是要求每条记录的损失数值更小，而是改变正负梯度比例。按110条记录平均，正系数为$(10\times1.1+90/90+10)/110=0.2$，负系数仍为0.8。

Bi-DEFUSE另把目标拆为窗口内转化概率$h_{\mathrm{in}}$与窗口外但$W$内的转化概率$h_{\mathrm{out}}$，以$h=h_{\mathrm{in}}+h_{\mathrm{out}}$合成。前者在首次等待后已得到完整二元标签，后者用延迟修正风险；共享层、任务专家和门控共同学习两类模式。两个事件互斥才可相加，实现应检查输出和不超过1以及总CVR校准，不能将两个任意sigmoid输出未经约束地解释成完整概率分布。

与ES-DFM相比，记录量取决于同一回流制度，增加的是类型标记、假负例辅助估计和可选多任务分支。若每条记录模型计算为$C_f$，主训练仍约为$O(N(1+\bar d)C_f)$，双分布版本额外增加分支代价。渐近修正依赖辅助概率、记录类型及采样条件正确；真实购买未被追踪到、同一订单重复回流或归因规则改变均不属于单纯的时间迟到。

\paragraph{DEFER：成熟负例回流的延迟学习}
\label{sec:rs-advertising-defer}

只补正例会让后续训练反复看到已确认的成功，却很少重新看到已确认的失败。\term{DEFER}在回流流程中加入成熟真实负例，并使用与其重复制度一致的权重\scite{rs-a18}它改变的是训练记录的组成，不是把“等待10分钟没购买”重新命名为成熟负例。

在平衡重复的理论版本中，每个点击有首次记录和一份真实标签记录：即时正例为正、正；迟到正例为负、正；窗口内未转化为负、负。即时正例也计入平衡重复，否则不能直接使用每个点击出现两次的分布公式。窗口结束前已有转化的点击不必等待$W$才确认正标签，未转化则须等到$W$才确认该窗口的负标签。相比之下，ES-DFM只回流迟到正例。

两次记录使点击特征的相对出现次数保持一致，但首次假负例仍改变正负比例。沿用$h$与$d$，混合记录中正比例为$h-d/2$，负比例为$1-h+d/2$，相应权重为
\begin{equation}
  w_+^{\mathrm{DEFER}}=\frac{h}{h-d/2},\qquad
  w_-^{\mathrm{DEFER}}=\frac{1-h}{1-h+d/2}.
  \label{eq:rs-advertising-defer}
\end{equation}
训练把$h$换为当前CVR估计、$d$换为成熟数据训练的迟到概率，计算加权交叉熵，并停止权重路径上的梯度。模型不能通过改变自身权重而不是改善预测来降低损失。分母过小、预测$d>h$等不一致需要约束或稳定化，并在验证中量化其影响。

仍取100次点击：10个即时正例贡献20条正记录，10个迟到正例贡献10条负与10条正，80个成熟负例贡献160条负。总计200条，其中30正、170负。$w_+=4/3$、$w_-=16/17$，故加权正负比例分别是$(30\times4/3)/200=0.2$与$(170\times16/17)/200=0.8$。成熟负例提供确定标签，又通过平衡重复避免只增加迟到人群的相对出现次数，但没有凭空产生未展示广告的响应。

每个点击保存至可回流时点，理论版本训练量约为$2N$，缓存压力还取决于$W$与特征大小。长窗口业务若提前用较短窗口近似成熟负例，可以降低存储和回流时延，却重新引入假负例，不能仍称标签全部正确。原文讨论了这一工程变体及多个延迟窗口的辅助任务；采用时应报告近似窗口、残余迟到率和校准变化。DEFER不解决追踪缺失，也不保证过去数周的成熟记录代表今天的新广告。

\paragraph{DDFM：双估计器的延迟响应学习}
\label{sec:rs-advertising-ddfm}

即使回流标签准确，回流数据也来自较早的点击，且其组成与新记录不同。\term{DDFM}把新观察流和回流流分开建立CVR风险估计，再按权重共同更新同一个响应模型\scite{rs-a19}这里的“双估计器”指两套风险估计，不意味着线上一定要部署两个独立CVR模型。

原文的新流含IP、FN、RN，回流流含DP与成熟负记录DN，不回流IP。令$a=1-h+d$为初次负记录概率，则新流中可见负标签的真实负比例为$r=(1-h)/a$；回流流中DP比例为$d/a$、DN比例为$(1-h)/a$。用统一函数
\begin{equation}
  \begin{aligned}
    \psi(\widetilde y,u;\widehat h)
      &=\widetilde y\ell_+
       +(1-\widetilde y)[(1-u)\ell_++u\ell_-],\\
    L_{\mathrm{DDFM}}
      &=\alpha\mathbb{E}_{\mathrm{new}}\psi(\widetilde Y,\widehat r;\widehat h)
        +(1-\alpha)\mathbb{E}_{\mathrm{return}}
           \psi(\widetilde Y,\widehat a;\widehat h),
        \qquad 0\leq\alpha\leq1,
  \end{aligned}
  \label{eq:rs-advertising-ddfm}
\end{equation}
分别平均两个流，避免流量比偶然变化替代设计的$\alpha$。新流中$u=r$修正假负例；回流流中$u=a$修正缺少即时正例等造成的组成偏差。后者给成熟负记录附加正损失项，不是在否认其真实标签，而是在用所见流估计目标流的总体风险。

辅助学习共用嵌入，分别训练真实负例比例$\widehat r$和迟到正例概率$\widehat d$。第二个风险系数采用$\widehat a=(1-\widehat h)+\widehat d$，其中CVR部分停止梯度。与直接计算多个概率之比相比，这种组合减少了额外除法，却不保证估计始终自洽；应检查$0\leq\widehat a\leq1$，并记录采取的约束。$\alpha$在时间隔离验证集选择，部署输出仍是$\widehat h$。

教学组的新流有10正、90负，$\widehat r=8/9$时，新流正损失系数为$0.1+0.9(1-8/9)=0.2$。回流流有10正、80负，$\widehat a=0.9$时，正系数为$1/9+(8/9)\times0.1=0.2$；两套负系数都是0.8。因此在这些条件概率恰好正确的同质数据中，任意固定$\alpha$均得到同一理想风险。真实流中若旧广告和新广告的特征分布不同，两套风险分别在各自特征分布上正确，并不自动等于“当前全部机会”的同一个边际风险。

\begin{table*}[htbp]
  \centering
  \small
  \begin{tabularx}{\linewidth}{@{}lXXX@{}}
    \toprule
    方法 & 10分钟未转化 & 2小时转化 & 24小时仍未转化\\
    \midrule
    ES-DFM & 初次负记录，负权重修正 & 回流正记录，正权重修正 & 不额外回流负记录\\
    DEFUSE & 正负损失混合，推断FN比例 & 回流DP，与IP使用不同权重 & 用成熟标签学习辅助概率\\
    DEFER & 首次记录按混合流加权 & 真实正标签回流 & 真实负标签回流；平衡版本每点击两份\\
    DDFM & 新流风险使用$\widehat r$ & 回流流正记录 & 回流流负记录，风险使用$\widehat a$\\
    \bottomrule
  \end{tabularx}
  \caption{延迟方法对同一时间线的处理。2小时转化与24小时未转化是两个互斥分支，不应让同一点击同时进入这两类最终记录。}
  \label{tab:rs-advertising-delay-methods}
\end{table*}

若即时转化比例为$\bar h-\bar d$，DDFM记录量约为$N(2-\bar h+\bar d)$，再加两个辅助头的训练；两个流的队列和来源标记必须分别维护。表~\ref{tab:rs-advertising-delay-methods}也说明，任意更换回流规则后沿用原权重都可能破坏推导。应按点击时间划分训练与未来测试，测试等待相同$W$成熟，并按广告新旧、延迟长度和流来源检查偏差。离线全量重排成熟记录可以验标签，不足以模拟在线可见性。

\subsubsection{概率校准与价值转换}
\label{sec:rs-advertising-calibration}

学到响应后，还要检查它是否能以概率使用。第~\ref{sec:rs-score-calibration}~节给出共同校准接口，第~\ref{sec:rs-score-drift}~节讨论更新。广告中尤其需要按广告主、位置、场景、素材版本和计费事件检查：一个总体CTR准确的模型，可能对首页高估、对搜索页低估；在按概率乘出价排序时，两类误差不仅影响报表，还改变谁获得曝光。

\paragraph{Field-aware Calibration：字段校准}
\label{sec:rs-advertising-field-calibration}

仅用分数做单调变换，同分样本总会得到相同结果，无法修正这种分组差异。\term{Field-aware Calibration}论文中的\term{Neural Calibration}，把冻结基础模型的logit $l=f(\bm{x})$与字段信息一起送入后处理器\scite{rs-a05}基础模型先在训练集拟合，独立校准集保存$(\bm{x}_n,l_n,y_n)$，未参与二者拟合的未来测试集负责验收。

校准输出为
\begin{equation}
  \begin{aligned}
    \widehat p_{\mathrm{cal}}
       &=\sigma\!\left(\eta(l)+g(\bm{x})\right),\\
    \eta(l)&=u_k+
         \frac{l-a_k}{a_{k+1}-a_k}(u_{k+1}-u_k),
         \quad l\in[a_k,a_{k+1}).
  \end{aligned}
  \label{eq:rs-advertising-field-calibration}
\end{equation}
$a_k$是固定logit结点，$u_k$是待学习的结点高度，$g$是字段修正网络。分段线性映射称为Isotonic Line-Plot Scaling（ILPS）。原文通过$\lambda\sum_k\max(u_k-u_{k+1},0)$惩罚下降区间，与校准集二元交叉熵共同优化；有限惩罚系数不等于严格投影约束，训练后仍应核验$\eta$单调。字段项可以改变不同组的相对次序，所以整体校准器并不保证保持基础排序。

教学假设广告位甲、乙各1000次曝光，基础模型全部预测0.1，实际分别有50、150次点击。总体200/2000正好为0.1，按这一分数合并的校准误差为0；但字段误差
\[
  \operatorname{F\!-\!ECE}
   =\sum_g\frac{N_g}{N}
       \left|\overline y_g-\overline p_g\right|
   =\tfrac12|0.05-0.1|+\tfrac12|0.15-0.1|=0.05.
\]
若$\eta$在此处保持原logit $\log(0.1/0.9)\approx-2.1972$，字段修正分别为$-0.7472$与$0.4626$，sigmoid输出便是0.05与0.15。未来测试若保持这些条件，字段误差降至0；若实际比例变了，旧修正不能被当作已验证的新概率。

按例~\ref{ex:rs-advertising-events}的每点击价值20元，两组每曝光预测价值由同样2元变为1元与3元。这里改变的是CTR及曝光价值，点击后CVR并没有因为位置CTR校准而自动改变。若实际按CPC支付，单次点击支付仍由机制决定；不能拿校准后的期望曝光价值直接记作当次实际支出。

训练代价由校准样本数和小型字段网络决定；用二分定位$K$个区间的线上代价为$O(\log K+C_g)$，其中$C_g$是字段网络计算。结点外输入需要明确截断或外推规则。小广告主缺少校准样本时，细分越多不一定越可靠，可共享参数并报告置信区间；数据发生变化后，更新校准器也须保留基础模型版本及时间边界。校准修正观察概率，不识别广告因果增量。

最终价值转换仍遵循式~\eqref{eq:rs-advertising-value}。CPC下可以从$\widehat h$与转化条件金额形成点击价值，CPM比较使用每曝光价值乘1000，目标CPA出价则使用允许成本与预测CVR建立控制输入。分组校准、金额估计和归因窗口任一变化都会改变出价依据，应分别追踪，而不是把所有收入变化归给排序模型。

\section{广告竞价与投放控制}
\label{sec:rs-advertising-control}

价值输入建立之后，系统仍有三个嵌套的决定：当前机会愿意报多少价，整个周期怎样使用资金，多个广告组和渠道之间怎样分配共同预算。平台决定分配和支付，广告主或其代理决定出价，资源控制器则调整参与和资金使用。这些决定的反馈尺度不同：一次点击可能几秒后计费，一次转化可能几天后确认，而跨渠道预算通常不能随每次请求任意转移。

\subsection{分层投放控制}
\label{sec:rs-advertising-hierarchical-control}

组合层向周期层提供预算和目标，周期层向竞价层提供参与概率、出价乘子及可用额度；执行层返回真实获胜、计费和转化记录。周期层不能只看“本轮赢了多少”，还要看用了多少余额、承担多少在途义务，以及后续机会是否仍有资格。硬频控继续接入第~\ref{sec:rs-presentation-time}~节的曝光状态，不作为提高出价即可放松的软目标。

\subsubsection{单次机会控制}
\label{sec:rs-advertising-single}

广告主面对的是一次有市场价格的购买机会。若以一次曝光为单位，记预测价值为$\widehat v$、最高有效竞争门槛为随机变量$M$，其条件分布为$F(m\mid\bm{x})$。单位置一价拍卖中，忽略并列概率和保留价时，出价$b$的获胜概率是$F(b\mid\bm{x})$，期望支出是$bF(b\mid\bm{x})$；准线性效用最大化求解
\begin{equation}
  b^*=\argmax_{b\geq0}(\widehat v-b)F(b\mid\bm{x}).
  \label{eq:rs-advertising-first-price-bid}
\end{equation}
若密度$f$存在且最优解在内部，则$(\widehat v-b)f(b)=F(b)$。报价因此不仅依赖响应价值，也依赖竞争分布。一价中的压价与预算节奏中的降价是不同原因：前者权衡中标概率和获胜后支付，后者还给后续机会保留资源。

教学假设$M$在0至10元均匀分布、当前曝光价值为8元，则效用为$(8-b)b/10$，最优$b=4$元，获胜概率0.4，期望支出1.6元、期望效用1.6元。相同竞争分布下，二价支出为$\mathbb{E}[M\mathbf{1}\{M\leq b\}]$；报价8元得到期望价值6.4元、期望支出3.2元和效用3.2元。两种支付规则不能共用“中标就扣出价”的模拟器。

实际\term{实时竞价}（RTB）中，$F$须按场景、时段和可观测竞争条件估计。完整门槛日志可以拟合条件分布；只观察输赢及获胜支付时，失败意味着门槛超过当前报价，属于删失信息，而不是一条价格为0的记录。按既有数据分布学习的$F$还依赖竞争者与底价是否稳定。可用离散价格网格比较候选$b$，计算代价约为网格数乘分布求值代价；分布尾部没有覆盖时，数值优化得出的高价未必有证据支持。

目标CPA提供另一种价值尺度。若允许每次转化成本$c$元，每点击预测转化概率为$\widehat h$，则$c\widehat h$元/点击是回报约束归一化后的价值输入。它不是单次点击已实现的利润。第~\ref{sec:rs-advertising-retrieval}~节原框架的优化CPC由预测CVR、商品价格及广告主营销比例相乘形成，同样要注明价格是收入、营销比例代表可用于广告的份额。多个点击共享成本约束时，周期控制再调整这个基础输入。

平台接收报价后，必须在明确的可行分配中确定获胜者和支付。来源、位置、混排与营销感沿用第~\ref{sec:rs-content-organization}~节及第~\ref{sec:rs-position-allocation}~节；Virtual Bids沿用第~\ref{sec:rs-virtual-bids}~节的机会成本接口。广告支付不是排序分数的自然单位，更不能把排序后因体验规则替换的广告继续按替换前的次序收费。

\paragraph{单位置一价与二价机制}
\label{sec:rs-advertising-single-auction}

设合格集合为$C$，报价$b_i$与\term{保留价}$r_0$都以元/同一种计费事件表示。这里采用先剔除$b_i<r_0$的广告、再取最高出价者的规则；并列按事先固定的广告标识顺序处理，未获胜者不支付。最高者为$i^*$，次高有效报价为$b_{(2)}$；没有次高者时，以保留价为竞争门槛。\term{一价}和\term{二价}事件支付分别是
\begin{equation}
  \tau_{i^*}^{\mathrm{first}}=b_{i^*},\qquad
  \tau_{i^*}^{\mathrm{second}}=\max(r_0,b_{(2)}).
  \label{eq:rs-advertising-single-auction}
\end{equation}
这些规则来自机制约定，不通过点击标签拟合。线性扫描$|C|$个报价即可取得前两名；资格查询、风控与并发预留另计。

教学假设甲、乙、丙真实单次价值为8、7、3元，提交报价6、4、2元，保留价1元，甲可用预算10元。两种机制都由甲获胜。一价支付6元、余额4元；二价支付4元、余额6元。若计费事件为点击且没有发生有效点击，两种机制的实际支出都是0，不能在曝光时直接把6或4元结为已支出；执行系统可以先预留相应额度。

单个独立物品、私人单参数价值、风险中性、效用为“获胜价值减支付”、没有跨拍卖预算或其他外部性、且保留价不随本人报价操纵时，二价真实报价是弱占优策略。给定他人最高门槛$m$，$m<v_i$时应赢、$m>v_i$时应输，报$v_i$恰好实现这一划分；报得更高不会降低支付，却可能赢下亏损机会，报得更低可能错过正效用机会。

一价没有这个结论，式~\eqref{eq:rs-advertising-first-price-bid}已说明真实价值和出价通常不同。即使支付规则是二价，当广告主最大化累计价值并受预算、CPA或回报约束时，也可能选择缩放价值报价；单次机制对准线性效用的诚实性，不等于自动竞价器应在每个请求直接报告未经调整的原始价值。

\paragraph{广义二价与质量调整}
\label{sec:rs-advertising-gsp}

多位置还需要声明点击如何随位置变化。采用可分离的教学设定：广告$i$在位置$j$的CTR为$q_i\alpha_j$，$q_i>0$是固定广告质量，$1\geq\alpha_1\geq\cdots\geq\alpha_K>0$是位置系数，与广告身份无关。广告提交每点击出价$b_i$，按$a_i=q_ib_i$排序；\term{GSP}的质量调整支付是下一名分数除以获胜者质量，另与每点击保留价比较：
\begin{equation}
  \tau_{i_j}^{\mathrm{GSP}}
   =\max\!\left(r_{i_j},\frac{a_{i_{j+1}}}{q_{i_j}}\right),
  \qquad
  \mathbb{E}[s_{i_j}]=q_{i_j}\alpha_j\tau_{i_j}^{\mathrm{GSP}}.
  \label{eq:rs-advertising-gsp}
\end{equation}
式中$i_j$为第$j$名，最后一个位置也需要落选竞争者或底价作为门槛。此处假设保留价过滤后的次序已固定。若实际使用分数底价、广告主差异底价或在排名后才检查底价，必须按相应规则重新推导。

\begin{table}[htbp]
  \centering
  \begin{tabular}{@{}lrrrr@{}}
    \toprule
    广告 & CPC出价$b_i$ & 质量$q_i$ & 分数$q_ib_i$ & 名次\\
    \midrule
    甲 & 6 & 0.10 & 0.60 & 2\\
    乙 & 5 & 0.20 & 1.00 & 1\\
    丙 & 3 & 0.15 & 0.45 & 3\\
    \bottomrule
  \end{tabular}
  \caption{两个位置的教学报价，金额单位为元；位置系数为1和0.5，每点击保留价均为1元。}
  \label{tab:rs-advertising-gsp}
\end{table}

表~\ref{tab:rs-advertising-gsp}中乙获得首位，每点击支付$0.60/0.20=3$元；甲获得次位，每点击支付$0.45/0.10=4.5$元。乙、甲每次展示的期望支出分别为0.6和0.225元。若这次仅甲被点击，则实际支出为甲4.5元、乙0元；期望支出不等于这次结算。质量来自响应预估及平台规则，而非广告主自报的支付价格。

GSP的“下一名定价”不自动带来真实报价。取质量均为1、位置系数1和0.5、真实每点击价值10、9、1元且无底价约束的另一教学例子。真实报价时第一名期望效用为$1(10-9)=1$；若它改报8元退到第二位，每点击支付1元，期望效用为$0.5(10-1)=4.5$。因此，它有获利偏离。固定质量和位置假设下研究出的均衡性质，不能改写为GSP的占优诚实性。

\term{VCG}使用不同支付原则：每名获胜者支付其占位给其他广告造成的总价值损失，而非只看下一名的单个报价。记可行分配集合为$A$，申报价值之和最大的分配为$a^*$，则$i$的期望支付为
\begin{equation}
  P_i^{\mathrm{VCG}}
   =\max_{a\in A_{-i}}\sum_{k\ne i}V_k^{\mathrm{bid}}(a)
      -\sum_{k\ne i}V_k^{\mathrm{bid}}(a^*),
  \label{eq:rs-advertising-vcg}
\end{equation}
其中$V_k^{\mathrm{bid}}$表示由申报价值构成的分配价值，不是系统预测器的输出；$A_{-i}$是没有$i$时的可行分配集合。实作应将申报值与模型估计分开存储。固定可分离CTR、无底价时，位置$j$的质量调整VCG每点击价格可进一步写成
\[
  \tau_{i_j}^{\mathrm{VCG}}
   =\frac{\sum_{\ell=j}^{K}
       (\alpha_\ell-\alpha_{\ell+1})a_{i_{\ell+1}}}
          {q_{i_j}\alpha_j},
  \qquad \alpha_{K+1}=0.
\]
按表~\ref{tab:rs-advertising-gsp}，乙的外部性是$0.5\times0.60+0.5\times0.45=0.525$元/展示，每点击2.625元；甲每点击仍为4.5元。本例保留价1元不约束结果。

VCG的占优诚实性依赖私人价值、准线性效用、真实可行域上的精确福利最大化和对应的外部性支付。预算约束、广告间价值外部性、由广告主可操纵的质量估计、近似组合求解或事后混排均可能破坏所用证明条件。一般布局下VCG需对删除各获胜者后的分配重新求解；若只在$K$个可分离位置排序，可利用上述差分公式，代价主要为$O(|C|\log|C|)$排序和线性累加。不能为了省一次求解就沿用无约束GSP价格声称仍是VCG。

\paragraph{自动竞价条件下的保留价与加性调整}
\label{sec:rs-advertising-reserves}

价值最大化的广告主可以在若干机会中承受较低回报，用其他机会的高回报补足周期约束，市场行为因而不同于逐次效用最大化。关于此类自动竞价者的研究利用有误差界的价值信号设置保留价与加性调整，并比较福利和平台收入\scite{rs-a08}其回报约束把周期约定价值与支付比较，不把二者之差作为目标。

为看清机制，固定单位置、同质质量和每次展示计费，令$r_i$为保留价，$\zeta_i$为预先给定且不受本人当前报价操纵的\term{加性调整}。先按申报$b_i$核验底价，再按$b_i+\zeta_i$竞争的版本中，获胜者$i^*$的二价门槛是
\begin{equation}
  \tau_{i^*}
    =\max\!\left\{r_{i^*},
       \max_{k\ne i^*,\,k\ \text{合格}}(b_k+\zeta_k)
          -\zeta_{i^*}\right\}.
  \label{eq:rs-advertising-boost}
\end{equation}
加性量提高排名时，也须从该获胜者的门槛中扣除；否则分配与支付不再对应同一个临界报价。保留价则限定其必须支付的最低值。质量调整、位置系数或延后底价核验加入后，要采用相应的完整机制，而非在式中随意混用单位。

教学假设甲、乙、丙报价为4、3、2元，保留价均为1元；不加调整时甲获胜，支付3元。现在$\zeta_{\text{甲}}=0$、$\zeta_{\text{乙}}=2$、$\zeta_{\text{丙}}=0$，三者分数变为4、5、2，乙获胜，支付$\max(1,4-2)=2$元。若乙保留价改为2.5元，则支付2.5元；若提高至3.5元，乙申报未过底价，在这里采用的预过滤版本中不参与竞争。这个表只演示规则，不声称提高乙权重必然改善真实福利。

原文还研究先按分数分配、获胜者未过本人底价便留空的延后核验方式，两者在不合格广告是否影响排名上不同。其福利近似保证需要价值信号满足指定乘法误差界、广告主采取相应未被支配策略，以及位置价值与竞价格式等假设；底价信号要求位于真实价值的适当下界区间。对部分非诚实机制还限制统一乘子策略。线上模型的总体CTR校准并不证明每名广告主的私人价值都满足这些界，半合成市场结果也不是任意市场下的提升保证。

\subsubsection{投放周期控制}
\label{sec:rs-advertising-period}

一段投放周期的状态至少包括剩余预算$B_t$、剩余时间、已实现价值与支出、当前机会特征及在途反馈。动作可以是是否参与、参与概率或价值出价乘子。按概率参与是随机跳过部分合格机会，缩放出价是改变可承受的市场门槛；两者可能得到相同平均花费，却购买不同质量的流量。频控、商品状态及预算预留仍由执行层逐次核验。

匀速支出只是一种目标曲线，不是价值最优的定义。若下午出现更高价值机会，早晨用完预算会损失这些机会；若全天缺少合格且有价值的流量，为了完成花费目标强行提价也可能损害回报。下面先比较可解释的反馈控制，再讨论从完整轨迹学习控制的模型。

\paragraph{AdaptivePacing预算节奏}
\label{sec:rs-advertising-adaptive-pacing}

\term{AdaptivePacing}在预算受限的效用目标中为支出加入对偶价格。eBay的研究以广告主最大每点击报价作为价值代理，采用CPC结算，并允许广告主给定分时目标花费曲线\scite{rs-a10}这与后文“最大化价值且满足回报要求”不同：本方法的目标已包含价值减支付。

单位置、固定质量下，若本次每点击价值代理为$\widehat v_t$、事件门槛为$\tau_t$，预算对偶变量$\mu_t\geq0$使有效效用变成$\widehat v_t-(1+\mu_t)\tau_t$。愿意购买的门槛对应
\begin{equation}
  b_t=\frac{\widehat v_t}{1+\mu_t},\qquad
  \mu_{t+1}
    =\Pi_{[0,\mu_{\max}]}
       \left[\mu_t+\eta_t(s_t-g_t)\right],
  \label{eq:rs-advertising-adaptive-pacing}
\end{equation}
其中$g_t$为本时段目标支出，$\sum_tg_t=B$；$s_t$是该时段实际支出，$\eta_t$是步长，$\Pi$表示投影。预算花快了则提高$\mu$、降低后续出价，花慢了则相反。非负投影来自对偶可行域，有限上界为工程控制参数，不能无说明地当作原理论保证的一部分。

教学假设三个时段目标各10元，总预算30元，初值$\mu_1=0$、步长0.1、上界不约束，当前每点击价值代理始终为6元。第一段实际花14元，$\mu_2=0.4$，后续基础出价由6降至$6/1.4\approx4.286$元；第二段实际花8元，$\mu_3=0.2$，基础出价回升至5元；第三段再花8元，$\mu_4=0$。累计实际支出为14、22、30元。三段支出是给定的教学反馈，不是仅凭这三个出价即可推导出来的市场结果。

每段只需常数个状态量及一次更新，在线乘子应用为$O(1)$；控制时段越细，对结算延迟和随机点击越敏感。初值可取0或经历史周期热启动，步长、目标曲线和可选上界须在过去或验证时段选择。流量曲线与价值曲线可以形成不同的目标支出，缺少机会时不应借降低底层质量门槛强行补齐花费。若业务增加最低参与、出价边界等变体，需要重新检查其对偶解释和最终约束，而不能把多个规则的局部效果相加。

剩余预算与在途预留是独立硬检查，乘子没有原子扣费能力。原文理论讨论、其多位置CPC部署以及业务变体也不能混为同一个证明对象；此处仅在给定单位置效用设定下推导出价。执行记录以机制确定的真实事件支付计算$s_t$，绝不能把中标报价代作二价实际花费。

\paragraph{顺序预算与回报控制}
\label{sec:rs-advertising-sequential-pacing}

如果预算服务和回报服务分别维护，容易先由回报控制器抬高或压低报价，再交给预算服务缩放。关于多约束节奏控制的研究把这种顺序结构与最小值结构、联合对偶结构放在同一目标下比较\scite{rs-a09}为保持口径，先把目标回报$\rho_{\mathrm{tar}}$归一化：若机会预测价值为$\widehat v_t$，则用$\bar v_t=\widehat v_t/\rho_{\mathrm{tar}}$，累计约束写成$\sum_ts_t\leq\sum_tu_t$，其中$u_t=\bar v_t a_t$，$a_t$是实际分配量。以下教学拍卖中$a_t$为获胜0或1，支付即时发生。

预算对偶变量为$\mu_t>0$，回报对偶变量为$\lambda_t>0$。单独的回报控制给出乘子$k_t^{\mathrm{R}}=(1+\lambda_t)/\lambda_t$，单独的预算控制给出$k_t^{\mathrm{B}}=1/\mu_t$。顺序结构把二者相乘：
\begin{equation}
  \widetilde b_t=k_t^{\mathrm{R}}\bar v_t,\qquad
  b_t^{\mathrm{seq}}
    =\min\!\left\{k_t^{\mathrm{B}}\widetilde b_t,B_t\right\}.
  \label{eq:rs-advertising-sequential}
\end{equation}
这里$1/\mu_t$在$\mu_t<1$时会放大而非降低报价；若业务只允许降价，会再截断该乘子，并形成需要另行分析的变体。

原文采用乘法反馈更新。用本章符号表达，其基本形式为
\begin{equation}
  \lambda_{t+1}
   =\lambda_t\exp[\eta_{\mathrm{R}}(s_t-u_t)],\qquad
  \mu_{t+1}
   =\mu_t\exp[\eta_{\mathrm{B}}(s_t-g_t)].
  \label{eq:rs-advertising-two-dual-update}
\end{equation}
回报控制器看最终执行后的价值与实际支出，预算控制器看相同实际支出与目标$g_t$，不能让前者误以为中间报价$\widetilde b_t$已经执行。乘法更新需要正初值；若把一个变量初始化为0，它会永远保持0，工程实现须避免这种吸收状态。真实转化迟到时，$u_t$若使用预测而非已实现转化，控制保证也只能先针对相应代理，后续还须校正成熟回报。

下面固定一条共同教学轨迹，供三种结构复算：三次机会归一化价值为$(2,1,3)$元，二价门槛为$(2,1.1,1.5)$元，总预算4元，$g_t=4/3$元；$\lambda_1=\mu_1=1$，两步长均为$\log2$，报价等于门槛时获胜。首轮顺序报价为4元，支付2元、取得价值2元，更新得$\lambda_2=1$、$\mu_2=2^{2/3}\approx1.5874$。次轮有效乘子为$2/1.5874\approx1.2599$，赢下价值1元、价格1.1元的机会。

此后$\lambda_3=2^{0.1}\approx1.0718$，$\mu_3=2^{13/30}\approx1.3503$。第三轮虽有较高的未截断报价，却只剩0.9元，买不起门槛1.5元。最终价值3元、支出3.1元，回报约束违反0.1元。预算服务确实未超支，回报服务也确实收到了反馈，二者顺序组合仍没满足整体回报要求。

\paragraph{最小值节奏控制}
\label{sec:rs-advertising-min-pacing}

最小值结构让两个控制器对同一基础价值独立报价，执行较小者：
\begin{equation}
  b_t^{\mathrm{min}}
   =\min\!\left\{
       \frac{1+\lambda_t}{\lambda_t}\bar v_t,
       \frac{1}{\mu_t}\bar v_t,
       B_t\right\}.
  \label{eq:rs-advertising-min-pacing}
\end{equation}
它仍用式~\eqref{eq:rs-advertising-two-dual-update}的真实反馈更新两变量，只在合成报价时共享最少的信息。不能把未获采用的控制器暂停更新，否则其状态不再反映整段实际执行。

在共同轨迹中，首轮回报乘子2、预算乘子1，由预算约束报价2元并获胜。次轮预算乘子降为$2^{-2/3}\approx0.6300$，执行报价0.6300元，放弃1.1元门槛；因为本轮$s_2=u_2=0$，$\lambda_3$仍为1，$\mu_3=2^{-2/3}\approx0.6300$。第三轮预算乘子回升至1.5874，仍小于回报乘子2，未截断报价约4.7622元，经余额截断为2元，支付1.5元并取得价值3元。

最终支出3.5元、价值5元，剩余0.5元并不意味着控制失败，因为没有更多可买机会。当前活跃约束也会切换：若另一个状态有$\lambda=2,\mu=0.4$，回报乘子1.5小于预算乘子2.5，就由回报约束报价。两种状态中的出价切换是有限协调，而不是重新联合求解全部未来机会。

每个请求维护两个标量、两次指数更新及一次取小，额外计算为$O(1)$。其保证不是由“取小一定安全”推出：每轮低报价仍可能累计违反回报目标，尤其当转化预测偏高时。上述研究在独立同分布机会、给定价值、诚实拍卖响应及相关正则条件下分析其遗憾和违反量；这不等于任意非平稳市场上的逐轮硬回报保证。

\paragraph{联合对偶约束优化}
\label{sec:rs-advertising-joint-pacing}

联合结构直接对同一周期的两个约束定价。目标和拉格朗日函数为
\begin{equation}
  \begin{aligned}
    &\max_{\{b_t\}}\sum_t\bar v_ta_t(b_t),\qquad
      \sum_ts_t(b_t)\leq B,\quad
      \sum_ts_t(b_t)\leq\sum_t\bar v_ta_t(b_t),\\
    \mathcal{L}
      &=\mu B+\sum_t
        \left[(1+\lambda)\bar v_ta_t(b_t)
          -(\mu+\lambda)s_t(b_t)\right].
  \end{aligned}
  \label{eq:rs-advertising-joint-objective}
\end{equation}
预算变量$\mu$只增加成本系数；回报变量$\lambda$同时提高取得价值的收益系数与支付的成本系数。这解释了为什么它们不能简单取平均。通用约束推导沿用第~\ref{sec:rs-goal-constrained-objectives}~节，应用关键在于价值和支付均按实际机制定义。

若单次机制对准线性效用是诚实的，令有效价值为$(1+\lambda_t)\bar v_t/(\mu_t+\lambda_t)$即可实现该单次拉格朗日最大化，因此
\begin{equation}
  b_t^{\mathrm{joint}}
    =\min\!\left\{
       \frac{1+\lambda_t}{\mu_t+\lambda_t}\bar v_t,
       B_t\right\}.
  \label{eq:rs-advertising-joint-pacing}
\end{equation}
若是一价，则应把这个有效价值送入式~\eqref{eq:rs-advertising-first-price-bid}对应的单次优化器，不能直接称该乘子报价是一价最优解。若$\mu_t+\lambda_t$接近0，需由明确的最高报价与剩余可支付额度限制，不能产生无穷报价。

在共同轨迹中，联合首轮乘子为1，与最小值结构一样支付2元。次轮乘子$2/(1+1.5874)\approx0.7730$，不同于最小值的0.6300，但同样未中标；第三轮乘子$2/(1+0.6300)\approx1.2270$，乘价值3后仍由2元余额截断，支付1.5元。其终局与最小值相同，途中报价并不相同。

\begin{table*}[htbp]
  \centering
  \small
  \begin{tabular}{@{}llrrrrrr@{}}
    \toprule
    结构 & 时段 & $\lambda_t$ & $\mu_t$ & 执行$b_t$ & 实际$s_t$ & 累计$S_t$ & 累计价值\\
    \midrule
    顺序 & 1 & 1 & 1 & 4 & 2 & 2 & 2\\
    顺序 & 2 & 1 & 1.5874 & 1.2599 & 1.1 & 3.1 & 3\\
    顺序 & 3 & 1.0718 & 1.3503 & 0.9 & 0 & 3.1 & 3\\
    最小值 & 1 & 1 & 1 & 2 & 2 & 2 & 2\\
    最小值 & 2 & 1 & 1.5874 & 0.6300 & 0 & 2 & 2\\
    最小值 & 3 & 1 & 0.6300 & 2 & 1.5 & 3.5 & 5\\
    联合 & 1 & 1 & 1 & 2 & 2 & 2 & 2\\
    联合 & 2 & 1 & 1.5874 & 0.7730 & 0 & 2 & 2\\
    联合 & 3 & 1 & 0.6300 & 2 & 1.5 & 3.5 & 5\\
    \bottomrule
  \end{tabular}
  \caption{三种控制结构的同轨迹教学比较。金额均为元；未截断报价经剩余预算限制后才执行。末列是已获归一化价值，不是平台收入，也不是观测到的因果增量。}
  \label{tab:rs-advertising-pacing-comparison}
\end{table*}

表~\ref{tab:rs-advertising-pacing-comparison}是检验接口的有限算例，不能证明长期优劣。上述多约束节奏研究将有限时域$T$下的价值与满足预算和回报约束的离线最优比较；在其独立同分布、支出与回报响应的连续性、局部稳定性及有界二阶矩等条件下，最小值和联合结构可达到$O(\sqrt T)$量级的遗憾与回报违反界，而顺序结构存在随$T$线性增长的违反或遗憾实例。小离散教学表不满足所有连续性假设，因此不以它验证该渐近定理。

联合结构每轮也是常数级对偶运算，但需要两个状态在同一时间口径下协调。模型预估价值、平台支付延迟、预算共享及市场反应可能成为主要误差来源。无论采用哪种结构，平均约束都不能替代逐次预算拒绝、频次核验与在途资金预留；离线使用完整最终转化进行控制，不能冒充真实流式策略。

\paragraph{RLB：价值计算的周期竞价}
\label{sec:rs-advertising-rlb}

对偶乘子概括了资金的稀缺程度，\term{RLB}则直接计算多花一单位预算会减少多少后续收益。它把剩余拍卖次数、预算和当前机会作为状态，依据竞争价格分布构造状态转移，求取累计价值，再由当前奖励与未来价值差决定出价\scite{rs-a20}原文以预测点击数为奖励；把奖励改为预测成交毛利是应用扩展，不是原实验已经验证的同一目标。

设还剩$n$次机会，剩余离散预算为$B$，当前特征为$\bm{x}$，CTR估计为$q(\bm{x})$。市场门槛$m$以同一最小货币单位离散化，条件概率质量为$p(m\mid\bm{x})$。记$\overline J(n,B)$为尚未观察下一次特征时的最优未来点击数，$J(n,B,\bm{x})$为已知当前特征时的价值。在二价、未来机会独立同分布等模型假设下，
\begin{equation}
  \begin{aligned}
    J(n,B,\bm{x})
     &=\overline J(n-1,B)\\
     &\quad+\max_{0\leq b\leq B}
       \sum_{m\leq b}p(m\mid\bm{x})
       \left[q(\bm{x})+\overline J(n-1,B-m)
                       -\overline J(n-1,B)\right],\\
    \overline J(n,B)&=\mathbb{E}_{\bm{X}}J(n,B,\bm{X}),
       \qquad \overline J(0,B)=0.
  \end{aligned}
  \label{eq:rs-advertising-rlb}
\end{equation}
方括号内是买入当前机会的净点击价值：取得$q(\bm{x})$，但失去花费$m$后本可用于未来的价值。余额减的是市场支付$m$，不是报价$b$。若广告位免费，零预算仍可能赢取免费机会，边界条件也应相应保留，不能一律令$\overline J(n,0)=0$。

市场门槛模型和CTR模型先从有相应观测的数据估计，动态规划再从$n=0$向较长周期计算。原文为简化计算使用全局或分段的价格分布近似，并以平均CTR构造低维状态值表；这种简化不应被误读为价格天然与用户特征独立。在线得到当前$q(\bm{x})$后，用未来值表评估不同价格阈值，选择累计期望收益最大的报价。

教学假设只剩两次机会、预算3元，当前门槛确定为3元、CTR为0.1；未来机会门槛确定为2元、CTR为0.3。当前买入后余额0元，累计期望点击为0.1；放弃后保留3元，未来可买，期望点击为0.3。因此当前3元支付的机会成本是0.3次点击，高于当前0.1，不应购买。若未来CTR只有0.05，比较方向就反转。策略依赖对未来的模型，不是“越接近结束就必然提高出价”的固定规则。

完整值表随时间和离散预算扩大。若最大价格网格有$M$点，用前缀累加评价候选报价，直接动态规划代价约为$O(TB_{\mathrm{grid}}M)$，保存全部时刻值表为$O(TB_{\mathrm{grid}})$；把价格上界当常数时才可隐藏$M$。原文的神经近似重点拟合相邻预算的价值差
\[
  \Delta J(n,B)=\overline J(n,B+1)-\overline J(n,B),
\]
以小规模动态规划得到的$(n,B,\Delta J)$为监督，用回归损失训练MLP。原因是实际出价依赖未来价值差，拟合很大的累计价值再相减，较小的相对误差也可能淹没单次点击收益。

部署时可由差分网络累加所需的预算机会成本，或通过分段周期与状态映射复用较小范围的近似。网络训练数据是计算出的价值标签，不能将其当作直接观测的长期真实收益。预算离散过粗会抹去临界价格，过细会使值表不可承受；新竞争者、价格分布漂移、CTR失准或压缩状态遗漏未来信息都会改变策略。第~\ref{chap:rs-long-term}~章的状态与策略评价条件在这里仍然适用。

\paragraph{DRLB：反馈学习的周期竞价}
\label{sec:rs-advertising-drlb}

RLB需要明确的价格与转移模型。\term{DRLB}改为每隔一个控制时段调整报价除数，用DQN从投放反馈学习调整动作，并单独学习能够体现整个周期收益的RewardNet\scite{rs-a21}它的动作频率低于单次请求频率：一次控制更新影响接下来一组曝光的报价。

状态$\bm{\xi}_t$包含当前时段、余额、剩余调整次数、最近预算变化率、已赢曝光的CPM、获胜率和上一时段取得的价值。令当前出价除数为$\kappa_{t-1}>0$，离散动作$a$关联调整比例$\beta_a>-1$，则
\[
  \kappa_t=\kappa_{t-1}(1+\beta_a),\qquad
  b_{t,i}=\frac{\widehat v_{t,i}}{\kappa_t}.
\]
这里$\kappa$不是支付价格，也不是前面联合对偶变量的简单改名；正调整提高除数、降低后续出价。每个请求仍要经过余额和资格检查，时段末汇总真实支出及价值形成下一状态。

仅用当段取得的点击奖励，早期提高报价容易显得有利，却可能提前耗尽预算。DRLB对一个曾访问的状态动作对，用包含它的已完成周期中的最大总收益形成训练奖励：
\begin{equation}
  \widetilde r(\bm{\xi},a)
   =\max_{e\in E(\bm{\xi},a)}
      \sum_{t=1}^{T}r_{e,t}.
  \label{eq:rs-advertising-drlb-reward}
\end{equation}
$E(\bm{\xi},a)$是包含该状态动作的历史周期集合。周期结束后更新字典中的最大值，RewardNet通过平方损失拟合这些标签。DQN再以RewardNet输出替代即时奖励，使用目标网络和经验回放拟合$\widetilde r+\max_{a'}Q_{\mathrm{target}}(\bm{\xi}',a')$；终止状态不再加入后继值。训练中原文还根据动作价值曲线是否符合单峰直觉调整探索强度，不把固定衰减的探索率当成唯一选择。

沿RLB的两机会例子，教学假设当前CTR价值为0.1、除数$\kappa=0.033$，未调整出价约3.03元，足以花完3元预算。若控制器选择增加8\%的除数，出价降至约2.81元，当前落选，未来仍可用预算购买2元门槛、CTR为0.3的机会。这个结果来自给定机会轨迹，是否能由日志学出该动作还取决于覆盖。再假设某状态动作曾出现在总收益4和5的周期中，RewardNet标签是5而不是4.5；另一个动作只出现在收益3的周期中，标签为3。这是对较好轨迹的乐观学习，不是两个动作的无偏因果比较。

原文关于新旧奖励共享最优策略的论证采用单一初始状态、确定性转移与固定周期长度等条件；随机市场上取样本最大值会受偶然好运、访问次数和策略覆盖影响，不能直接套用一般奖励塑形的不变性保证。RewardNet和Q模型也可能相互强化错误，因此评价仍需回到原始累计价值、总支出与约束违反，而非只看训练奖励提高。

设每次DQN、RewardNet前向代价为$C_Q,C_R$，候选动作数为$|A|$，每个控制时段的动作选择约为$O(|A|C_Q)$，每个曝光应用除数仅为$O(1)$；训练另需两个网络及经验池存储。状态汇总遗漏未来机会、预算记录延迟和奖励代理变化都会影响学习。固定日志回放假设未执行机会的价格和响应可获得，若新策略改变竞争者出价，就不再是同一环境。

\paragraph{DiffBid（AIGB）：状态轨迹生成的周期竞价}
\label{sec:rs-advertising-diffbid}

将每个状态单独压缩后做价值自举，可能遗漏长历史中影响后续机会的信息。\term{DiffBid}把预算等状态组成整条轨迹，以条件扩散学习轨迹与收益、约束条件的关系，再从生成的下一状态恢复竞价参数；\term{AIGB}是其AI-Generated Bidding范式的名称\scite{rs-a22}这不意味着所有广告任务在任何状态表示下都违反马尔可夫条件，也不意味着生成模型的未来状态已经真实发生。

离线样本是一段日志$\omega=\{(\bm{\xi}_t,\bm{a}_t,r_t)\}_{t=1}^{T}$。状态含剩余时间、余额、花费速度和当段及累计成本效率等，动作$\bm{a}_t$为竞价参数向量。将状态堆成$\bm{X}^{(0)}\in\mathbb{R}^{T\times d_s}$，条件$\bm{c}$包含由该完整训练轨迹计算的归一化收益、成本约束是否达标等属性。训练时可以用完整轨迹建立监督；推断时只能固定已经观察到的状态前缀，不能读取当天最终回报作为输入事实。

对随机扩散步$k$给状态轨迹加噪，噪声网络$\epsilon_\theta$学习在条件下去噪，逆动力学网络$f_\phi$从一段历史状态与下一状态预测日志动作。两项训练目标为
\begin{equation}
  \begin{aligned}
    L_{\mathrm{DiffBid}}
      &=\mathbb{E}\left[
          \left\|\bm{\epsilon}
            -\epsilon_\theta(\bm{X}^{(k)},\bm{c},k)\right\|_{\mathrm{F}}^2
        \right]\\
      &\quad+\mathbb{E}\left[
         \left\|\bm{a}_t-
           f_\phi(\bm{\xi}_{t-L:t},\bm{\xi}_{t+1})
         \right\|_2^2\right].
  \end{aligned}
  \label{eq:rs-advertising-diffbid}
\end{equation}
$L$是可见历史长度，$\|\cdot\|_{\mathrm{F}}$为矩阵Frobenius范数。训练随机丢弃部分条件，以支持有条件与无条件噪声预测的组合引导；推断从噪声轨迹出发，反复固定真实历史、去噪未来，取下一预测状态交给逆动力学。通用生成建模见\BookChapterRef{chap:pgm-generative}{模型卷的“概率生成模型”章}，本节重点是状态生成和动作执行之间的接口。

沿三时段预算4元的教学轨迹，设当前真实余额4元，条件要求较高累计价值且回报达标。生成器可能提出余额序列$(4,2,2,0.5)$，对应预测支出$(2,0,1.5)$；逆动力学给出三个价值乘子$(1,0.6,1.5)$。在前述门槛$(2,1.1,1.5)$、价值$(2,1,3)$下，这些参数恰好实现买入、放弃、买入，实际支出3.5元。这里的乘子是教学假设的网络输出，并非从余额序列唯一可解：许多报价都可能赢得相同支付，同样的报价在不同市场也可能产生不同余额。

真实执行后只把实际支付写回余额，再重新规划。若首轮市场门槛变成2.5元，报价2元会落选，真实余额仍为4元；不能强制把状态改成生成的2元余额以“保持轨迹一致”。约束条件只能引导生成分布，硬预算、频控、最高报价和在途额度仍在执行端裁决。

设一次完整轨迹噪声网络计算为$C_\epsilon(T,d_s)$、逆动力学为$C_f(L)$，单个去噪样本训练约为两者之和；采用$K_{\mathrm{diff}}$步去噪的在线规划约为$O(K_{\mathrm{diff}}C_\epsilon+C_f)$，尚未计入多个候选轨迹。网络若对时间做全注意，$C_\epsilon$本身包含随轨迹长度增长的代价，不能只报“线性于去噪步数”而省略它。规划可在较慢控制时段执行，单次曝光只应用已确定参数。

离线日志中没有的高收益条件、过强引导、逆动力学不可辨识及市场漂移均可能产生不可实现轨迹。DiffBid原研究报告了其特定数据和平台下的离线及在线比较；这些证据支持相应设定中的性能，不证明生成出的任意轨迹都满足预算和成本要求。训练条件与部署目标相同，也不等于已识别广告主的真实增量价值。

\paragraph{GAVE：价值引导的生成竞价探索}
\label{sec:rs-advertising-gave}

只模仿日志中的行动，模型难以改进已有策略；任意扰动行动，又可能离开日志能够支持的范围。\term{GAVE}以带目标条件的Decision Transformer预测竞价乘子，在训练时围绕日志动作产生探索标签，用预测的后续得分比较两种标签，并学习一个偏向较高收益的序列价值作为探索方向\scite{rs-a23}这里的探索主要发生在离线训练中，不表示已经把新动作交给真实市场试投。

样本仍是完整投放轨迹，输入交错排列的收益到终点、状态和动作。与只累计转化不同，GAVE先把成本要求并入轨迹得分。设截至$t$累计转化代理为$U_t$、支出为$S_t$，目标CPA为$c$，一个常用得分及其收益到终点为
\begin{equation}
  \begin{aligned}
    F_t&=U_t\min\left\{1,
         \left(\frac{c}{S_t/U_t}\right)^\zeta\right\},\\
    R_t^{\mathrm{go}}&=F_T-F_{t-1}.
  \end{aligned}
  \label{eq:rs-advertising-gave-score}
\end{equation}
$\zeta>0$控制超成本惩罚；无转化时需预先定义零得分，无支出而有合法免费转化时惩罚系数取1。若$U_t$来自预测概率之和，式中成本是每单位预测转化的成本，不是成熟实际CPA。分数也不一定随时间增加，因此$R_t^{\mathrm{go}}$不必非负。更换业务目标时要重新计算训练条件及评价分数，而不是只替换推断提示中的几个词。

记日志动作$a_t$为正的价值乘子。Transformer同时输出待学习动作$\widehat a_t$、探索系数$\widehat\beta_t$、后续得分估计和序列价值。探索系数由$\sigma(u_t)+0.5$给出，故$\widetilde a_t=\widehat\beta_ta_t$处于日志动作的0.5到1.5倍之间。将原动作和探索动作分别接到同一历史后，得到预测后续得分$\widehat R_{t+1}$和$\widetilde R_{t+1}$；比较权重为
\[
  w_t=\sigma\!\left(\alpha_R
       (\widetilde R_{t+1}-\widehat R_{t+1})\right).
\]
动作损失把$\widehat a_t$同时拉向两种标签：
\[
  L_a=\mathbb{E}\left[
      (1-w_t)(\widehat a_t-a_t)^2+
      w_t(\widehat a_t-\widetilde a_t)^2\right].
\]
更新这项损失时冻结$w_t$和$\widetilde a_t$的梯度，避免模型通过改变标签来降低误差。得分预测头另用真实日志中的$R_{t+1}^{\mathrm{go}}$作平方误差监督；探索动作的高分则是模型估计，二者不能混为同等可靠的标签。

为了使探索有方向，价值头以非对称平方损失拟合较高的条件期望分位（expectile）：
\[
  L_V=\mathbb{E}\!\left[
    \left|\eta-\mathbb{I}\{R_{t+1}^{\mathrm{go}}-\widehat V_{t+1}<0\}\right|
    (R_{t+1}^{\mathrm{go}}-\widehat V_{t+1})^2\right].
\]
原文取接近1的$\eta$，强调较好的日志回报，再使探索动作的预测得分接近冻结的$\widehat V_{t+1}$。这个统计量不是分位数，也不是有限样本下得到认证的最大可达收益上界。完整训练联合得分拟合、动作拟合、非对称价值拟合和探索得分引导，各损失权重由独立验证周期选择。

教学假设日志动作$a_t=1$，探索系数1.2，原动作与探索动作的预测后续得分分别为4和4.4，$\alpha_R=1$，则$w_t\approx0.599$。只考虑本样本的动作损失，其最小点为$(1-w_t)\times1+w_t\times1.2\approx1.120$，而不是直接跳到1.2。另设周期已有支出8元、预测转化2次，目标CPA为3元，$\zeta=2$，则得分为$2(3/4)^2=1.125$；多取得转化但成本失控的轨迹，可能得分更低。两处数字都是教学假设，不是论文性能。

线上只需将真实历史和预设目标交给动作头，取$\widehat a_t$，再以$b_{t,i}=\widehat a_t\widehat v_{t,i}$形成当前时段报价。训练中的探索标签不直接作为线上动作。设历史长度为$L$、隐藏维度为$d$，全注意主干单层主要计算为$O(L^2d+Ld^2)$；训练还需要对原动作及探索动作评价，线上则主要承担一次动作预测。乘子应用仍是逐请求常数运算。

限制相对扰动只能控制动作距离，不能保证拍卖结果距离：出价从0.99变到1.01，跨过1元门槛就会从不花费变成花费。高期望分位可能放大日志中的偶然好运，得分头在未见动作上也可能过于乐观。GAVE的模拟及特定部署比较应与日志覆盖、成本违反和市场设定一起阅读；硬预算仍由执行端保证，不能由价值引导替代。

\paragraph{GUIDE：行为备选与价值比较的竞价}
\label{sec:rs-advertising-guide}

探索动作有可能更好，也可能只是价值模型的误判。\term{GUIDE}保留两个动作来源：Decision Transformer直接预测动作与下一状态，逆动力学模型根据当前状态和预测下一状态重建更接近日志行为的动作，再由双Q网络比较两者\scite{rs-a24}这是一种备选比较结构，不是先后执行两次竞价。

训练样本包含$(\bm{\xi}_t,a_t,r_t,\bm{\xi}_{t+1},a_{t+1})$及其历史、收益到终点条件。DT预测$(\widehat a_t,\widehat{\bm{\xi}}_{t+1})$，逆动力学给出
\[
  \widehat a_t^{\mathrm{idm}}
      =f_\phi(\bm{\xi}_t,\widehat{\bm{\xi}}_{t+1}).
\]
动作与状态平方误差监督DT，日志动作$a_t$监督逆动力学。训练初期在预测下一状态处停止梯度，使逆动力学的误差不干扰尚未稳定的DT；随后开放这条梯度路径，让动作重建也约束状态预测。状态、金额与比率须按训练统计缩放，否则一个大量级余额误差可能淹没动作误差。通用Transformer机制见\BookChapterRef{chap:transformer}{模型卷的“Transformer”章}。

两个Q网络从日志转移学习，目标使用目标网络的较小值：
\begin{equation}
  y_t=r_t+\delta(1-d_t)
       \min_{k\in\{1,2\}}
          Q_k^{-}(\bm{\xi}_{t+1},a_{t+1}),\qquad
  L_Q=\mathbb{E}\sum_{k=1}^{2}
          (Q_k(\bm{\xi}_t,a_t)-y_t)^2.
  \label{eq:rs-advertising-guide-q}
\end{equation}
$d_t$是周期终止指示量，$\delta$是折扣，$a_{t+1}$为采样日志中的后续动作，不能擅自换成全动作空间的最大值。目标网络按指数移动平均更新。DT的目标再减去生成动作的$\min(Q_1,Q_2)$，鼓励高估计价值行动；逆动力学保留行为重建职责，不把这项探索目标直接加给它。这样才保留探索与行为备选的区别。

推断时分别计算
\[
  \underline Q_{\mathrm{dt}}=\min_k Q_k(\bm{\xi}_t,\widehat a_t),
  \qquad
  \underline Q_{\mathrm{idm}}=\min_k Q_k(\bm{\xi}_t,\widehat a_t^{\mathrm{idm}}),
\]
选择较大的一个所对应的动作。教学假设DT给出乘子1.2，其双Q预测为$(5.2,4.6)$；逆动力学给出0.9，其双Q预测为$(4.9,4.8)$。前者取小值4.6，后者取4.8，因此选择0.9。当前预测价值2元、二价门槛1.7元、余额3元时，出价1.8元获胜，实际支出1.7元、余额1.3元。这里4.8是未来累计价值估计，1.7才是当前支付，不能相减来更新资金。

如果下一状态预测不可能实现，逆动力学仍可能给出一个数值动作；低重建误差并不构成物理可达性的充分条件。同样，两个Q共享日志偏差时，取最小值仍可能共同高估。拟合旧策略后继价值再选择新动作，也需要单独验证新策略的累计效果，不能把Q排序当成无偏政策评价。

设DT前向代价为$C_{\mathrm{DT}}(L)$、逆动力学为$C_f$、单个Q为$C_Q$，每次控制选择约为$C_{\mathrm{DT}}+C_f+4C_Q$，另有双目标网络和训练存储。原文在AuctionNet及其规定的竞争者设置下检验模型；行为备选不是任意市场下的安全策略。所有候选动作都须经过同一额度、价格和资格核验，并记录最终执行动作，才能识别Q选择与安全截断各自产生的影响。

\paragraph{AIGB-Pearl：轨迹评价与竞价策略搜索}
\label{sec:rs-advertising-pearl}

单步Q需要从后继估计中自举，误差可能沿周期传播。\term{AIGB-Pearl}为完整轨迹训练独立评价器，再固定评价器优化生成规划器；Pearl在原文中展开为Planning with EvaluAtor via RL\scite{rs-a25}评价器直接拟合轨迹质量，不以自己的下一时刻预测构造标签，但这只去掉一种误差来源，并未消除离线分布外误差。

样本为状态轨迹$\omega$、其真实记录的质量$y(\omega)$和广告主信息。数值轨迹由序列网络编码，商品标题、类别、价格及历史表现等可用文本经语言模型编码后补充广告主表示。所谓专家信息在这里是可追溯的广告主上下文和预训练表示，不是让语言模型凭空给当天收益打分；历史统计只能使用规划前的信息。

评价器$\widehat y_\phi(\omega)$采用三类监督。逐轨迹平方误差拟合分数的大小；从日志中选出优劣轨迹对$(\omega^+,\omega^-)$，用
\[
  -\log\sigma\!\left(
       \widehat y_\phi(\omega^+)-\widehat y_\phi(\omega^-)\right)
\]
学习相对顺序；另对相近轨迹的分数差加入Lipschitz惩罚，抑制过度敏感。例如惩罚
$[|\widehat y_\phi(\omega_1)-\widehat y_\phi(\omega_2)|
  -L_e\|\omega_1-\omega_2\|_{\mathrm F}]_+$，
其中$[u]_+=\max(u,0)$。最小化的成对损失必须带负号，才能鼓励优轨迹得分高。只有成对监督时，分数的绝对尺度未被固定，不能直接当作收入金额。

规划器按目标质量$y^\star$和真实历史，自回归生成下一状态的高斯分布。训练先拟合日志条件分布，再在冻结评价器下提高所生成轨迹的得分，同时保留条件行为克隆及条件敏感性限制：
\begin{equation}
  \begin{aligned}
    L_{\mathrm{plan}}
      ={}&-\mathbb{E}_{\omega\sim p_\theta(\cdot\mid y^\star)}
                  \widehat y_\phi(\omega)
       -\lambda_{\mathrm{bc}}\,
          \mathbb{E}_{(\omega,y)\sim D}\log p_\theta(\omega\mid y)\\
       &+\lambda_{\mathrm{lip}}\,
          \mathbb{E}\left[
          \widehat W_1(y_1,y_2;\theta)-L_p|y_1-y_2|
          \right]_+ .
  \end{aligned}
  \label{eq:rs-advertising-pearl}
\end{equation}
这里$\widehat W_1$用同一组高斯噪声分别生成两个条件下的轨迹，再计算其距离，作为Wasserstein距离的耦合上界。共享噪声有助于把差异归于条件变化，而不是两次无关抽样。有限样本惩罚只是约束的实现近似，不能写成训练后已经严格满足全空间Lipschitz条件。

教学假设两条轨迹的记录质量分别为5和3，评价器预测为4.6和3.4，则平均平方误差为0.16，成对损失为$-\log\sigma(1.2)\approx0.263$。若规划器在两者上的概率为$(0.4,0.6)$，预测期望分数是3.88；改为$(0.6,0.4)$后为4.12。不计正则时，分数上升0.24支持这次搜索方向，但一个从未执行过、被误评为8分的异常轨迹更可能吸引优化。行为克隆和敏感性惩罚正是为了限制这种利用评价误差的现象，仍须检验实际轨迹质量。

推断沿用真实状态前缀，生成下一状态，经动作恢复接口输出当前竞价参数，执行后再接入实际反馈。完整未来用于训练和搜索，不写入真实余额。若一次序列生成需$T$步、单步代价为$C_p$，评价器代价为$C_e(T)$，每轮$M$条搜索样本约需$O(M(TC_p+C_e))$；训练还为配对条件生成额外轨迹。部署只规划所需下一步时可减少开销，但时延须按实际实现测量。

原文的评价偏差界依赖日志拟合误差、轨迹质量及评价器的平滑性、生成分布与日志的距离，以及目标条件的外推范围。不能删去这些条件，只保留“有理论保证”一句。第~\ref{chap:rs-evaluation}~章区分的代理评价与真实效果在此尤其重要：评价器分数提高是搜索依据，模拟或平台实验才检验投放效果，而预算不超支来自执行停止与额度检查，并非评价器自己证明。

\paragraph{AIGB-R1：规划与执行的分层竞价}
\label{sec:rs-advertising-r1}

宏观判断需要结合广告主偏好和较长历史，逐请求报价则要求确定的数值接口与低时延。\term{AIGB-R1}把两者分开：语言模型在投放周期前形成结构化策略提示，轻量Prompt Decision Transformer（\term{PDT}）结合提示与数值历史输出竞价乘子；随后通过模拟交互共同调整两层策略\scite{rs-a26}这里采用2026年7月作者预印本的研究设定，不把模拟后训练称为真实市场在线学习。

策略提示包含枚举标签和文字说明，例如风险偏好、预算节奏及其历史依据。PDT将这些提示嵌入为前缀，再接入历史收益到终点、状态和动作。离线预训练先让语言模型为已有轨迹标注策略，再用日志行动最大化条件似然。为了支持连续动作探索，执行器输出正值乘子的对数正态分布：
\[
  \log a_t\mid p,\bm{\xi}_{t-L:t},a_{t-L:t-1},R_{t-L:t}^{\mathrm{go}}
       \sim\mathcal N(\mu_{\theta,t},\sigma_{\theta,t}^2).
\]
训练最小化动作负对数似然。精确的零动作不属于对数正态支持，停止参与须由单独资格规则或明确的数据处理承担，不能把零直接取对数。离线轨迹标注可利用完整训练轨迹，部署提示却只能来自当前已知信息。

后训练对同一广告主任务生成$M$个策略提示，每个提示在模拟器中执行$K$条轨迹，得到得分$F_{m,k}$。模拟器按实际成交门槛扣费、返回转化及成本，以带CPA惩罚的周期分数作反馈。每项任务保存历史较好策略和最近尝试，用于下一轮生成；历史最优是已尝试集合中的最优，不是市场全局最优。

两层参数不能不加区别地分享同一个随机收益。令$\overline F_m$是提示$m$下的平均得分，$\overline F$是全部得分均值，原文的Decoupled Group Relative Policy Optimization（\term{D-GRPO}）利用分解
\begin{equation}
  F_{m,k}-\overline F
    =(\overline F_m-\overline F)
     +(F_{m,k}-\overline F_m).
  \label{eq:rs-advertising-r1-credit}
\end{equation}
组间部分评价规划提示，组内部分评价同一提示下的执行差异；分别归一化后，用对应策略的概率比及截断目标更新规划器与执行器。这样减少一次偶然执行结果对宏观提示的误记账，但它是学习中的信用分解，不是广告转化归因。

教学假设两个提示各模拟两次，得分矩阵为
$\left(\begin{smallmatrix}4&6\\3&5\end{smallmatrix}\right)$。
两提示均值为5和4，总均值4.5，因此规划器的未归一化信号为$(0.5,-0.5)$；两个提示内执行器的信号均为$(-1,1)$。第一提示第二条轨迹的总偏差1.5，恰为$0.5+1$。若每个提示仅执行一次，组内偏差必为零，无法用这个分解比较同一提示下执行器的好坏。

沿三个时段预算4元的例子，规划提示可以是“保守起步，保留后段预算”，执行器基于逐段历史给出乘子$(0.9,0.7,1.1)$。在价值$(2,1,3)$和二价门槛$(2,1.1,1.5)$下，前两段落选，第三段支付1.5元，余额2.5元。这个假设轨迹既可能是保留资金的结果，也可能是投放不足；仅凭语言理由不能判断是否合理，须与相同机会流中的可行替代策略比较。

一次后训练任务需$MK$个完整模拟周期，加上规划器生成、PDT执行和两层参数更新，成本随候选提示、重复次数及周期长度共同增加。部署把语言规划放在慢路径，PDT放在控制时段，逐曝光仅应用有效乘子；提示解析失败、超时或脱离可执行标签时，需要预定的降级策略。模拟器错误、离线标注泄漏、语言建议与数值执行不一致以及新市场竞争都限制迁移。结构化提示提高可检查性，却不能证明预算、CPA或真实增量达标。

\begin{table*}[htbp]
  \centering
  \small
  \begin{tabularx}{\linewidth}{@{}lXXX@{}}
    \toprule
    方法 & 状态与输出 & 学习或更新依据 & 比较时必须保留的条件\\
    \midrule
    AdaptivePacing & 花费误差、预算；报价除数 & 实际支出推动对偶更新 & 目标曲线、步长、在途义务\\
    顺序控制 & 两个乘子；乘积报价 & 共用最终执行反馈 & 局部控制可能形成整体违反\\
    最小值控制 & 两个乘子；较紧报价 & 分别更新预算与回报状态 & 约束切换与机会分布\\
    联合对偶 & 两个对偶量；联合报价 & 同一拉格朗日目标 & 价值口径、稳定性假设\\
    RLB & 剩余机会、资金；阈值报价 & 价格模型、动态规划值差 & 价格删失、网格与转移误差\\
    DRLB & 时段状态；除数调整 & 周期最大收益标签、Q学习 & 乐观奖励、日志覆盖\\
    DiffBid & 状态前缀；生成轨迹及逆动作 & 去噪与动作重建监督 & 未来可达性、生成计算\\
    GAVE & 目标及历史；价值乘子 & 探索标签、得分与价值引导 & 得分代理、分布外价值误差\\
    GUIDE & 历史；直接及逆动力学动作 & 行为重建、双Q比较 & 两种备选的共同估计偏差\\
    AIGB-Pearl & 目标及轨迹；生成规划 & 独立轨迹评价与受限搜索 & 评价器外推、约束近似\\
    AIGB-R1 & 策略提示及历史；连续乘子 & 离线预训练、分层模拟后训练 & 模拟器、提示与执行一致性\\
    \bottomrule
  \end{tabularx}
  \caption{周期算法的共同执行接口与不同证据。比较时固定预算、机会流、竞争者、计费方式和奖励口径，分别报告累计价值、实际支出与约束违反；表中没有将不同论文的收益数字横向拼接。}
  \label{tab:rs-advertising-cycle-methods}
\end{table*}

\subsubsection{广告组合与跨渠道控制}
\label{sec:rs-advertising-channels}

单个广告组能否买到一个机会，由其报价、资格及竞争决定；广告主还要决定哪些组或渠道获得多少资金。渠道可能只接受每日预算和目标回报，不开放逐次出价。因此，跨渠道控制的动作应写成渠道接口参数，不能假设广告主能够调用渠道内部每场拍卖的最优行动。

设渠道$j$收到预算$u_j$和目标价值回报$\rho_j$，在当前机会集合$\mathcal E_j$上返回价值$V_j(u_j,\rho_j;\mathcal E_j)$和实际支出$D_j(u_j,\rho_j;\mathcal E_j)$。共同预算$B$和总体回报$\rho$约束的是
\[
  \sum_j\mathbb{E}[D_j]\leq B,\qquad
  \sum_j\mathbb{E}[V_j]\geq\rho\sum_j\mathbb{E}[D_j].
\]
分别要求每个渠道回报不低于$\rho$，可能排除整体可行的组合；反过来，各渠道分别声称完成本地目标，也不能证明共享资金没有重复占用。渠道返回的归因价值还可能重复记入同一订单，必须先用共同的时间、货币、转化窗口和信用制度建立可加账目。

跨渠道预算研究表明，在渠道按其约束最优采购、机会与成本分布不随接口参数改变、允许分数分配等设定下，仅优化各渠道预算可以达到直接选择全部机会的基准，而仅调整渠道回报目标可能任意差\scite{rs-a12}这是接口能力的条件结论，不是“目标回报在实际平台上没有用”。真实渠道存在学习期、最低预算、不可拆订单、缓存和竞争反馈时，需要重新检查响应集合。

以两个渠道为例，一个渠道拥有少量高回报机会，另一个拥有大量回报略低但仍有价值的机会。统一提高两个渠道的目标回报，可能同时砍掉第二渠道的大部分规模；给第二渠道一个有限额度，则能明确控制消耗，让第一渠道贡献的回报余量支持组合。究竟应分多少钱仍未知，需要取得不同预算水平下的响应证据，而不能从当前单点平均ROAS直接推断边际收益。

\paragraph{SGD--UCB渠道预算学习}
\label{sec:rs-advertising-sgd-ucb}

\term{SGD--UCB}把连续渠道预算离散为可试验的档位，用上置信界探索未知的渠道响应，再用对偶随机梯度协调总预算和回报。它只观察本轮真正选中的档位所产生的价值与花费，不假设同时知道所有未选预算的结果。

为说明原算法，设每期平均预算额度为$\bar B$，共学习$T$期；渠道本地回报目标设为0，整体回报$\rho$仍保留。原文进一步假设每个渠道在允许预算范围内能花完分配额度，因而$D_j(u)=u$，且每个可能的机会实现都有严格满足总体回报的正机会。预算档位集合为$A=\{0,\Delta,\ldots,\bar B\}$，末档不超过$\bar B$。记录每个渠道档位的访问数$n_{j,t}(u)$和价值均值$\overline V_{j,t}(u)$。

若价值已归一化到$[0,1]$，原算法采用形如
$r_{j,t}(u)=\sqrt{2\log T/n_{j,t}(u)}$的置信半径；其他已知有界尺度要相应缩放。先访问各档，避免零分母，随后选择
\begin{equation}
  u_{j,t}\in\arg\max_{u\in A}
    \left\{\overline V_{j,t}(u)+r_{j,t}(u)
        -\frac{\lambda_t\rho+\mu_t}{1+\lambda_t}u\right\}.
  \label{eq:rs-advertising-channel-choice}
\end{equation}
$\lambda_t$对应总体回报，$\mu_t$对应平均预算。由拉格朗日式
$(1+\lambda)V_j(u)-(\lambda\rho+\mu)u$可见，置信上界鼓励试验，最后一项给资金定价。它不是简单按当前平均回报从高到低分钱。

渠道反馈后，只更新已选档位的计数和样本均值，两个对偶量使用
\begin{equation}
  \begin{aligned}
    \lambda_{t+1}
      &=\Pi_{[0,C]}\!\left(
         \lambda_t+\eta\sum_j(\rho u_{j,t}-V_{j,t})\right),\\
    \mu_{t+1}
      &=\Pi_{[0,C]}\!\left(
         \mu_t+\eta\left(\sum_j u_{j,t}-\bar B\right)\right).
  \end{aligned}
  \label{eq:rs-advertising-channel-dual}
\end{equation}
投影上界$C$、步长、网格、保底小预算和总期数共同构成算法参数。原算法还维护累计回报余量及累计支出，用剩余期数检查是否仍有足够余量继续探索；触发停止条件后，转入预设的小预算方案。这个方案的可行性依赖前述严格可行假设，不能在任意渠道上凭空指定一个“小额投放必然赚钱”的后备策略。

\begin{example}[两渠道的一次预算学习]
\label{ex:rs-advertising-channel}
教学假设$T=128$，两个非零预算档为0.5和1，零预算的价值已知为0。下表给出已积累的反馈摘要；价值和资金均为归一化单位。令$\lambda_t=0.5$、$\mu_t=0.8$、$\rho=0.8$，则式~\eqref{eq:rs-advertising-channel-choice}的单位资金系数为0.8。
\begin{center}
  \begin{tabular}{@{}ccrrrr@{}}
    \toprule
    渠道 & 预算 & 均值 & 访问数 & 置信半径 & 选择分数\\
    \midrule
    甲 & 0.5 & 0.70 & 32 & 0.551 & 0.851\\
    甲 & 1.0 & 0.90 & 32 & 0.551 & 0.651\\
    乙 & 0.5 & 0.40 & 32 & 0.551 & 0.551\\
    乙 & 1.0 & 0.55 & 8  & 1.101 & 0.851\\
    \bottomrule
  \end{tabular}
\end{center}
甲选择0.5，乙选择1；乙的大预算获选是因为不确定性更大，而非均值回报更好。假设两个渠道均花完额度，返回价值0.65和0.60，则实际总支出1.5、总价值1.25，回报余量为$1.25-0.8\times1.5=0.05$。若$\bar B=1.25$、$\eta=0.2$且投影未起作用，则$\lambda_{t+1}=0.49$、$\mu_{t+1}=0.85$。甲0.5档的新均值为$(32\times0.70+0.65)/33\approx0.6985$，乙1档为$(8\times0.55+0.60)/9\approx0.5556$。

本期比平均额度多花0.25，只能由学习总周期的已有余量支持。若业务要求每期最多1.25，则本次独立档位选择不可直接执行，必须增加共同额度分配层；修改后的算法不再原样享有原定理。置信上界超过1也不表示预测概率超过1，它是这里未截断的探索评分。
\end{example}

最终输出可取各期渠道预算的平均配置。原文的分段线性、递增、凹响应来自渠道内的分数采购模型；在该结构、独立稳定抽样及预算耗尽等条件下，适当选择网格和步长，可得到随学习期数收敛的平均配置误差及总体约束结论。若$J$个渠道各有$K$档，直接扫描每期为$O(JK)$，均值和访问计数存储也是$O(JK)$，真正昂贵的成本是取得各档位反馈所需的投放期数和试验资金。

当渠道不能花完预算时，式~\eqref{eq:rs-advertising-channel-dual}不能继续把分配额当实际支出，选择目标中的成本响应也要另行学习。若提高预算会改变竞价者策略、响应曲线随季节漂移，旧档位均值和置信半径也不再自动有效。更根本地，若$V_j$是归因收入，本节学得的是该核算目标下的配置；只有具备后文的增量证据，才能解释为增量价值配置。噪声归因统计、聚合模型和实验校准都应把不确定性传给预算比较。

\subsection{投放学习与评价}
\label{sec:rs-advertising-policy-evaluation}

单次出价、周期控制和跨渠道配置有不同的比较单位。一次机会上的更高预期收益，不代表周期结束后更好；一个渠道增加收入，也可能来自另一个渠道或自然流量的替代。评价须从要改变的策略范围出发，保存该范围内的状态及竞争条件。

\subsubsection{投放日志与策略识别}
\label{sec:rs-advertising-policy-logs}

广告主可见的日志常比平台日志更窄。落选时可能只知道竞争门槛高于自己的报价；一价获胜只知道门槛不高于报价，不会从自己的支付中观察到最高对手价。预算耗尽后的请求甚至不再进入投放方日志。因而“把新报价与旧成交价比较”只有在竞争门槛记录完整、相应位置机制一致等条件下才构成合理回放。

日志应分开记录机会到达、资格、是否参与、请求时余额与预留、原始动作、安全处理后的动作、分配、计费事件、实际支付、响应标签及成熟时间。旧策略的行动概率也应在决策时保存，不能事后用新模型回填。若只保存获胜曝光，就无法区分没有机会、主动放弃和竞价失败；若只保存最终日预算，就无法恢复中途的资金资格。

例如，旧策略上午花完100元，下午没有参与记录。新策略上午少花30元后究竟能在下午取得什么，既需要下午机会与竞争信息，也需要对应条件下的响应证据。把旧日志的下午收益填成零，会把旧策略停止观测误当成环境没有价值。反之，直接借用另一广告主的下午成交，也不能保证用户、资格和价格可比。

第~\ref{chap:rs-evaluation}~章的倾向加权或双重稳健估计可以用于有明确支持的策略比较，但不能创造未出现状态的证据。周期策略需要处理整条行动历史如何改变状态；逐机会独立加权不足以恢复不同预算路径。连续出价密度比可能极端，确定性旧控制器也可能没有新动作支持，应报告覆盖、权重诊断和不能识别的范围。

延迟转化另有成熟问题。新策略若较晚花钱，实验结束时转化可能比早花钱的旧策略更不完整；不统一结果窗口就会把支付节奏差异当成转化率差异。评价响应模型时检验成熟标签上的校准，评价投放策略时则检验同一结算协议下的周期价值与支出。训练用的延迟损失修正不能自动替代实验的结果等待或删失处理。

\subsubsection{市场反馈与在线评价}
\label{sec:rs-advertising-market}

提高一个竞价者的出价，可能改变其他竞价者的获胜与支付；平台扩大广告负载，也可能改变未来用户访问和供给。市场不是固定的价格表。在线评价至少同时观察平台实际收入、广告主约定价值、预算与回报违反，以及第~\ref{chap:rs-presentation}~章的负载、营销感、频次和拒绝反馈，不能将其中一项代理为全部收益。

预算和回报约束还改变福利的定义。关于竞价器互动的研究使用\term{受支付能力限制的福利}（liquid welfare）：给定各广告主最终获得的价值$V_i$、预算$B_i$和目标价值回报$\rho_i$，
\begin{equation}
  W_{\mathrm{liq}}=\sum_i\min\{B_i,V_i/\rho_i\}.
  \label{eq:rs-advertising-liquid}
\end{equation}
它度量在这些约束下各方最多愿意支付的总量，不等于实际平台收入，也不等于无预算截断的总价值\scite{rs-a11}教学假设两广告主的$(B_i,V_i,\rho_i)$分别为$(4,12,2)$和$(5,6,2)$，则该福利为$4+3=7$；即使实际支付共5元，两者仍是不同指标。

该研究在价值序列按固定分布独立抽取、所有竞价者执行其指定算法、拍卖满足相应分配与支付性质等条件下，给出相对于最优事前受限福利的约一半保证，另有随周期增长次线性的误差项。论证不要求竞价动态先收敛到均衡，但不能据此断言任意学习器互相竞价都有效率。个体遗憾使用满足逐轮平均约束的动态乘子基准，并受市场漂移路径长度和响应平滑性影响；它不是任意对抗环境下相对于所有全周期策略的无遗憾保证。

实验随机化单位必须覆盖主要干扰路径。同一广告主的两组若共享预算，处理组更快消耗会减少对照组的投放资格；两组广告在同一拍卖互相竞争，又会改变对照价格。用户级随机化能隔离用户处理，但未必隔离竞价和资金。可按广告主或市场簇分组、分配独立额度，或使用带清洗期的时间切换；每种设计都要明确估计的是局部试验效果还是全面采用后的市场效果。

投放后的供给响应还可能更慢。广告主看到短期回报下降，可能降低后续预算、修改素材或退出，平台收入随之变化。第~\ref{chap:rs-long-term}~章讨论的行为与供给反馈在广告市场中同时存在。短期实验与较长周期观察应使用各自结论，不能用固定竞争者的一天回放替代广告主长期参与效果。

\paragraph{AuctionNet：多竞价者交互评价基准}
\label{sec:rs-advertising-auctionnet}

\term{AuctionNet}将机会生成、多个竞价代理和可配置拍卖连接起来，并提供生成日志与基线算法，以支持相同条件下的控制策略比较\scite{rs-a27}它是环境与数据，而不是接在响应模型后面的另一种出价预测器。

机会模块先用编码器和解码器学习历史机会特征的潜在表示，训练包含重建误差与潜在分布正则；在潜在空间训练噪声预测网络，再从去噪后的表示解码机会。价值模块结合生成特征、广告主特征和时间信息，以记录中的价值标签做平方误差监督，得到每个竞价者对机会的价值矩阵。其训练目标是接近数据分布及标签，不是证明生成机会具有正确的因果效果。通用生成和注意机制分别见\BookChapterRef{chap:pgm-generative}{模型卷的“概率生成模型”章}与\BookChapterRef{chap:transformer}{模型卷的“Transformer”章}。

每个控制时段，代理只取得允许观察的自身状态和机会信息，输出价值乘子；拍卖模块根据全部报价确定位置与价格，再用实际支出更新各代理余额。代理可使用节奏控制、线性规划、离线强化学习或生成策略。多位置、底价、计费单位和预算不足时的处理都应作为配置保存，不能一边按二价训练、一边按一价支付评价却不报告。

教学假设一次单位置二价机会对甲、乙的价值分别为2元和1元，乘子均为1，甲出2元、乙出1元，甲支付1元。若乙乘子变为1.5，获胜者仍为甲，甲却支付1.5元。甲预算3元时，两种竞争条件分别留下2元和1.5元，后续可行行动也不同。这说明冻结旧成交价的日志无法检验对手改变策略后的资金轨迹。

固定竞争者协议只替换待测代理，保持其他代理策略、初始预算、机会种子和机制不变，并轮换待测角色；它回答给定对手下的改进。共同更新协议允许多个代理一起学习，除相同初值和训练资源外，还需记录更新频率、价格轨迹、各方约束和退出行为；它回答学习互动是否稳定。两者不应把得分混合成一张不注明协议的排名表。

设有$N$个代理、每时段$M$次机会、$T$个时段。直接按报价排序的拍卖计算约为$O(TMN\log N)$，若只取少数获胜位置可采用相应选择算法；还要加上机会生成、$N$个代理决策及反馈汇总。保存逐机会逐代理记录的空间随$TMN$增长，轨迹汇总更小但会丢失门槛删失和位置细节。实验应报告这些成本和随机种子，避免让更多模拟资源被误认为算法本身更高效。

生成数据与真实数据的分布投影相近，只能检查部分统计性质，不能验证竞争行为、极端预算、季节变化或稀有转化。预生成日志和可交互环境也有不同信息条件。AuctionNet适合复现固定条件比较及竞争敏感性分析，真实共享预算下的增量和市场反馈仍需独立实验。

\section{广告创意优化}
\label{sec:rs-advertising-creative}

出价决定为机会付出多少，创意决定用户实际看到什么。同一背包可以突出电脑隔层、重量或价格，也可以选用不同图片和布局。生成任务提供有依据的素材候选，选择任务决定哪个版本获得当前展示与测试流量；如果把二者混在一起，就难以判断效果变化来自素材本身、受众变化还是探索分配。

\subsection{创意生成}
\label{sec:rs-advertising-creative-generation}

候选可以由已有标题、图片、背景和模板组合，也可以通过模型生成新文字或新图像。前者搜索离散组合，后者扩展素材内容，均需符合商品事实、价格、品牌规范与落地页。个性化只改变表达重点，不改变商品属性：用户喜欢登山，不会使一款防泼水通勤包自动获得防暴雨或专业承重能力。

第~\ref{sec:rs-presentation-expression}~节讨论对象怎样被表达，本节限定广告生成的输入和验收接口。输入应有可追溯事实及版本，输出要能分解出每项主张并找到证据。语言模型认为“更吸引人”只是一种生成或评价信号；审核、真实响应和长期体验分别承担不同验收职责。

\subsubsection{HLLM-Creator个性化标题生成}
\label{sec:rs-advertising-hllm-creator}

将全部点击历史直接拼到生成提示中，既昂贵，也难以让模型稳定利用兴趣。\term{HLLM-Creator}用Item LLM编码点击过的广告，再用User LLM聚合历史，最后由Creative LLM结合用户表示、原标题和卖点生成标题\scite{rs-a14}这里依据2025年作者预印本及其抖音搜索广告设定，通用文本生成机制不在本节重复展开。

训练样本包含历史点击序列$H_u$、目标广告事实$F_i$、可用查询及合成标题$Y_{u,i}$。合成数据先从历史提取兴趣描述，再选取与目标广告有关的兴趣组织标题，最后根据原标题和卖点过滤无依据内容。这样构造的$Y_{u,i}$是受事实检查的生成监督，不是用户亲自写出的偏好标签。原文使用语言模型过滤幻觉；部署仍应对价格、数值、资质等可核验字段采用权威数据核对，并抽查模型过滤器的漏检。

Item LLM将每项广告文字及末尾的专用标记编码，取该标记表示$\bm{e}_i$；User LLM接收按时间排列的$\bm{e}_{i_1},\ldots,\bm{e}_{i_n}$及用户标记，形成$\bm{u}$。Creative LLM把$\bm{u}$映射到自己的表示空间，作为目标广告事实词元之前的条件，逐词生成标题。一个历史广告只需编码成向量即可被多个历史引用，避免反复对整段历史文字做联合注意，但陈旧或未授权的历史仍应排除。

主损失是合成标题的负对数似然：
\[
  L_{\mathrm{gen}}
    =-\frac{1}{|Y|}\sum_{\ell=1}^{|Y|}
       \log p_\theta(Y_\ell\mid\bm{u},F_i,Y_{<\ell}).
\]
原文另加三个辅助目标。推荐损失用用户和候选广告表示预测点击，正例来自点击，负例随机采样；语义对齐用InfoNCE拉近同一用户的历史表示与兴趣描述表示，区分批内其他用户；重建损失由$\bm{u}$解码兴趣描述。三者分别要求表示能区分候选、接近文字兴趣、保留可解码的兴趣内容，总目标是$L_{\mathrm{gen}}+\lambda_{\mathrm{cls}}L_{\mathrm{cls}}+\lambda_{\mathrm{align}}L_{\mathrm{align}}+\lambda_{\mathrm{recon}}L_{\mathrm{recon}}$。

小批量梯度训练这些共享表示及生成参数，验证时分别看事实错误、标题质量和兴趣利用。随机采样负例下的点击辅助分数不自动是线上校准CTR；合成兴趣文本与点击历史也并非独立证据。因此，辅助目标降低不等于用户真实偏好被因果识别，更不等于广告增量提高。

教学假设同一商品已经核实“20升、重量650克、15英寸电脑隔层、防泼水面料、售价399元”。通勤组输入可强调电脑携带，候选标题为“399元轻量通勤包，15英寸电脑独立收纳”；有登山兴趣的用户组若仍匹配这件商品，可以生成“650克轻量日用背包，20升收纳”，但不能生成“专业登山包，暴雨全天防水”。前两者每个数值和属性都能回查，第三个既新增用途等级，又夸大防护。若投放当天售价变为429元，含399元的版本必须更新或撤下，不能以训练时事实正确为由继续使用。

为降低生成量，原文先对用户表示做$K$类聚类，再用广告与用户组中心的匹配预测器选择前$k$组，仅为这些组生成标题。教学假设有100万用户、1000个广告、$K=256$且$k=8$，逐用户逐广告会有10亿个潜在生成对，按组筛选后仅有8000个广告--用户组生成对；这不是端到端加速倍数，因为还需编码历史、聚类和计算匹配。已有标题随后由线上素材选择接口取用，而不是让每次曝光等待完整生成。

若每段历史含$n$项、每项文字长$\ell$，Item LLM对各项独立编码的全注意部分随$n\ell^2$增长，User LLM对$n$个表示的全注意部分随$n^2$增长；生成仍取决于目标标题长度和Creative LLM规模。聚类减小服务对象数，也会合并组内差异；匹配截断可能使少数需求没有个性化素材。模型裁判与同类过滤器还可能共享错误。真实效果应在同一投放和受众条件下比较，分别报告事实通过、点击、转化及负反馈，不能只用语言模型偏好分替代。

\subsection{创意选择}
\label{sec:rs-advertising-creative-selection}

事实合格的素材也未必表现相同。给每个版本均分流量，可能把大量测试消耗在明显较差的组合；总是选当前最高均值，则可能永远无法了解新版本。第~\ref{sec:rs-exploration-exposure}~节的探索接口在这里以素材版本为动作，学习所需成本不仅是曝光，还包括审核、制作、用户重复接触和放弃其他创意的机会。

素材轮换还应接入第~\ref{sec:rs-presentation-time}~节的素材、商品及品牌历史。一张新图不意味着用户第一次见到该广告主，换标题也不能重置硬频控。关闭、屏蔽和“不感兴趣”等交互条件见第~\ref{sec:rs-presentation-interaction}~节；短期点击更高的版本可能带来更多误导或疲劳，测试需同时观察这些结果。

\subsubsection{AutoCO创意组合与探索}
\label{sec:rs-advertising-autoco}

\term{AutoCO}把模板、图片、文字、背景等元素组合成创意特征，搜索元素之间适合的交互函数，用近似参数后验表示不确定性，再通过Thompson采样选择展示版本\scite{rs-a13}它优化已有素材的组合与试验分配，不负责生成新图像，也不以点击学习代替视觉合规检查。

一条样本$(\bm{x}_{p,c},y)$包含商品$p$、实际展示创意$c$的元素字段和二元点击$y$。假设有$L$个字段，每个字段值映射为嵌入。不同字段对可从拼接、逐元素乘、加、最大值和最小值五种操作中选择交互，并用全连接映射统一维度。若每对字段都枚举五种操作，结构数为$5^{L(L-1)/2}$；这一数量是结构搜索空间，不是线上每次一定计算这么多版本。

AutoCO为每个字段对维护连续操作权重，前向选择稀疏操作，按点击交叉熵的梯度更新连续权重，并投影回允许范围。每种操作使用自己的字段嵌入，以减少加法、乘法等不同梯度形态相互干扰。训练后按每对字段的最大操作权重形成交互结构。与第~\ref{sec:rs-fm}~节固定内积的FM相比，它改变的是交互函数及相应参数，不是把每种完整创意都作为互不相关的新ID。

记选定结构的点击概率为$\widehat q_{\bm{\theta}}(\bm{x})$，其中$\bm{\theta}\in\mathbb R^P$是需要采样的参数向量。对其设置高斯先验，以对角高斯$Q_\psi(\bm{\theta})$近似后验，最小化
\begin{equation}
  L_{\mathrm{VI}}
    =D_{\mathrm{KL}}(Q_\psi(\bm{\theta})\|P(\bm{\theta}))
      -\mathbb{E}_{\bm{\theta}\sim Q_\psi}
        \sum_{(\bm{x},y)\in D}
         \log\!\left[
           \widehat q_{\bm{\theta}}(\bm{x})^y
           (1-\widehat q_{\bm{\theta}}(\bm{x}))^{1-y}
         \right].
  \label{eq:rs-advertising-autoco}
\end{equation}
重参数化$\widetilde{\bm{\theta}}=\bm{\mu}+\bm{\sigma}\odot\bm{\epsilon}$，其中$\bm{\mu},\bm{\sigma},\bm{\epsilon}\in\mathbb R^P$，$\bm{\epsilon}\sim\mathcal N(\bm{0},\bm{I}_P)$，$\bm{I}_P$为$P$阶单位矩阵，即可用小批量梯度联合更新均值、正标准差和操作参数。该后验是变分近似，不是已校验覆盖率的置信区间。

一次请求先取得合格组合集合，采样一份参数$\widetilde{\bm{\theta}}$，用同一份采样比较所有候选，展示预测CTR最大的创意，随后把版本和点击反馈加入日志。共享参数意味着使用相同元素的版本会共同受益于学习，不能给每个版本独立抽一个不相容的“参数世界”却仍称为原来的参数Thompson采样。

\begin{example}[四个创意版本的一次选择]
两种模板甲、乙与两张已审核图片一、二组成四版。以下为教学假设，后验均值预测CTR依次为0.080、0.075、0.060、0.070，贪心策略会选择甲一。一次共同参数采样产生的CTR为0.073、0.084、0.078、0.081，故Thompson策略选择甲二。

为了检查这组抽样值可来自共同模型，令logit由偏置$-2.5$、模板乙效应0.2、图片二效应0.1及四项交互组成。取交互值约为$-0.0415,0.0108,-0.1698,-0.2288$，所得概率经舍入为上述抽样值。共享元素不要求所有交互同号，也不意味着四个CTR独立。

若甲二本次没有点击，仅增加$(\bm{x}_{p,\text{甲二}},0)$这一观测；不能给其余三版补零。点击标签也未说明成交，若后续目标改为毛利，还需接入相应价值和延迟反馈。一次探索损失或收益不能代表累计策略优劣。
\end{example}

若候选合格组合数为$K_c$，每版计算全部字段对、嵌入维度为$d$，忽略操作映射的常数后，直接比较代价约为$O(K_cL^2d)$；操作感知嵌入和变分均值、方差增加存储。候选组合过多时仍需约束枚举、缓存或分阶段截取，后验采样本身不解决组合爆炸。

新增和替换素材会改变候选集合，旧参数可共享已有元素的证据，但新元素仍需要探索。非平稳受众、曝光位置变化、后验方差过小或操作搜索过拟合都可能破坏探索效果。实验应固定商品、受众、位置和竞价条件，保留测试流量占比，比较累计点击或约定价值与探索代价。若创意改变点击质量、转化延迟或用户疲劳，还要扩大评价窗口；素材美观、当次响应与累计投放效果并不是同一目标。

\section{广告效果测量}
\label{sec:rs-advertising-measurement}

一次订单既可以进入广告主的销售报表，也可以被分配给某次广告触点，但这些记录不能回答没有广告时是否仍会购买。本节先区分两种测量目标：\term{转化信用归因}对已经发生的结果记账，增量识别比较不同投放条件下的结果。追踪受限改变可得数据及发布条件，却不取消这一区别。

\subsection{转化信用归因}
\label{sec:rs-advertising-attribution}

给定一笔已观察转化$k$及其窗口内的合格触点集合$P_k$，归因规则输出$\gamma_{k,i}\geq0$。若要求全额分配，则$\sum_{i\in P_k}\gamma_{k,i}=1$；没有可归触点的订单应进入“未归因”类别，不必强行分给最后一次可见广告。转化金额$M_k$分给渠道$j$的账目为
\[
  R_j^{\mathrm{attr}}
    =\sum_k M_k\sum_{i\in P_k}
       \gamma_{k,i}\mathbb I\{\operatorname{channel}(i)=j\}.
\]
每一笔金额只被分配一次时，各渠道之和才与原始转化金额守恒。跨平台各报一份全额归因，会让简单相加超过真实收入。

路径记录应包含触点的时间、素材、点击或浏览类型、渠道、设备关联依据及资格窗口。短窗口排除早期触点，跨设备缺失会让路径看起来更短，重复上报会让同一触点出现多次。规则和模型都只能处理允许且能够观察的路径，不能从缺失记录中证明用户没有接触广告。

\subsubsection{首触与末触信用核算}
\label{sec:rs-advertising-first-last}

\term{首触}将全部信用分给窗口内第一个合格触点，\term{末触}分给最后一个。规则参数来自核算制度，包括窗口长度、点击优先与否、浏览资格、去重及转化价值定义，无须从标签拟合。必须先确定这些参数再寻找“第一”或“最后”，不能根据希望得到的渠道表现临时改规则。

沿同一教学路径，假设转化发生前第6天有展示A、第2天有点击B、当天有点击C；三者都在7天窗口内，A属于展示渠道，B、C属于搜索渠道，一笔订单收入400元。若展示和点击均有归因资格，首触信用为$(1,0,0)$，展示渠道记400元；末触信用为$(0,0,1)$，搜索渠道记400元。订单并没有因规则变化多卖一单，只是信用位置改变。

若窗口改为3天，A被排除，首触改由B获得；若规定只有点击合格，即使保留7天窗口，A也不能取得首触。重复上报的C需要按事件标识去重。若真正末触D发生在无法关联的另一设备上，当前末触只能称为“可观察路径中的末触”，不应冒称完整用户旅程的最后触点。

在已按时间排序的路径上，首末触选择和渠道汇总为$O(|P_k|)$，若输入未排序则另有排序成本。规则简单可审计，但它既没有估计各触点的反事实结果，也没有刻画自然购买概率。将首触改成末触后某渠道ROAS上升，不能直接证明应给该渠道增加预算。

\subsubsection{DARNN双注意力归因}
\label{sec:rs-advertising-darnn}

固定规则不使用触点内容或用户行为。\term{DARNN}（Dual-attention Recurrent Neural Network）以展示特征序列和点击序列学习转化预测，再把展示与点击两组注意权重动态混合，形成触点信用\scite{rs-a15}通用循环网络见\BookChapterRef{chap:recurrent-networks}{模型卷的“循环神经网络”章}；这里的目标是完整观察路径的预测与核算，不是当前请求尚未发生时的CTR预估。

样本为$(\bm{x}_{1:m},c_{1:m},y)$，$\bm{x}_j$包含第$j$个触点的渠道、素材和情境，$c_j$是点击，$y$是约定结果窗口内是否转化。LSTM编码器顺序读取展示特征，形成$\bm{h}_j$；点击解码器结合前一点击、前一隐状态和编码器的最终状态$\bm{h}_m$形成$\bm{d}_j$，以二元交叉熵预测点击序列。由于解码器可读取完整展示路径的编码，不能把这种离线辅助预测当成无需未来信息的在线点击模型。

展示注意力与点击注意力分别通过能量网络和softmax形成
$\alpha_j^{\mathrm{imp}}$、$\alpha_j^{\mathrm{clk}}$，各自在整条路径上和为1。加权得到两个上下文向量
$\bm{c}^{\mathrm{imp}}=\sum_j\alpha_j^{\mathrm{imp}}\bm{h}_j$、
$\bm{c}^{\mathrm{clk}}=\sum_j\alpha_j^{\mathrm{clk}}\bm{d}_j$。
一个由末触特征和上下文计算的门控$\lambda\in[0,1]$，决定点击分支的权重：
\begin{equation}
  \widehat p
    =f_{\mathrm{conv}}\!\left(
       (1-\lambda)\bm{c}^{\mathrm{imp}}
          +\lambda\bm{c}^{\mathrm{clk}}\right),\qquad
  \gamma_j=(1-\lambda)\alpha_j^{\mathrm{imp}}
                  +\lambda\alpha_j^{\mathrm{clk}}.
  \label{eq:rs-advertising-darnn}
\end{equation}
$f_{\mathrm{conv}}$以sigmoid输出转化概率。点击分支注意的是点击模式的隐状态，并非只能给实际点击过的位置非零权重。

训练先用逐触点点击损失稳定序列表示，再加入最终转化的二元交叉熵联合训练。正转化和未转化路径都参与概率学习，路径窗口与成熟规则须一致。信用计算时，对已发生转化读出$\gamma_j$并汇总金额；模型训练并没有观察一组“正确因果贡献比例”来监督这些权重。

对上述三触点路径，教学假设展示权重为$(0.5,0.3,0.2)$、点击权重为$(0.1,0.4,0.5)$、$\lambda=0.6$，则混合信用为$(0.26,0.36,0.38)$，和为1。400元分别记104、144、152元，展示渠道104元、搜索渠道296元。若该路径的转化预测为0.2且实际转化，预测损失为$-\log0.2\approx1.609$；这项损失检验概率任务，并不证明上述分账比例是真实增量比例。

原文还以归因渠道回报配置预算并回放历史，对预算耗尽后缺失必要触点的路径取消转化计数。这是一种指定规则下的代理回放；用户可能在没有该触点时仍购买，也可能因新的触点组合改变行为，因此不能无条件把回放值视为真实策略收益的因果下界。第~\ref{chap:rs-evaluation}~章的预测评价和策略识别应分别执行。

若路径长度$m$、隐维度$d$，双LSTM主要代价约为$O(md^2)$，注意聚合为线性于路径长度，训练需保留相应隐状态。长路径截断、身份关联错误、末触信息优势及历史选择都会影响权重。循环隐状态已混入前面的触点，最后又经过非线性转化网络，因此注意权重不是可加的因果贡献，甚至不必等于删除一个触点后的模型预测变化；后者同样只是模型内部的敏感性分析。

\subsection{广告增量识别}
\label{sec:rs-advertising-incrementality}

增量测量先定义处理。令$A_u=1$表示用户被分配到某项广告投放政策，$A_u=0$表示规定的对照政策，$Y_u(a)$为相同结果窗口下的潜在转化、收入或其他结果。目标可以是
\[
  \Delta=\mathbb E[Y_u(1)-Y_u(0)],
\]
也可以是给定处理前特征的条件平均效应。这里的对照必须具体：不展示本广告、改展其他广告、展示公益广告或减少整个广告负载，会产生不同答案。

增量成本回报也要与该比较一致。若实验政策比对照多支出$\Delta S$，多带来$\Delta N$次转化和$\Delta R$收入，则增量CPA为$\Delta S/\Delta N$，增量ROAS为$\Delta R/\Delta S$，分母接近零时需特别报告不确定性。归因订单总量不能代替$\Delta N$，处理组全部销售额也不能代替$\Delta R$。

高响应与高增量并不相同。重定向已把商品加入购物车的用户，可能获得很高成交率，但其中很多订单没有广告也会发生。相反，响应率较低的人群可能拥有更大的可改变空间。增量选择需要估计两种条件之间的差，而不是给CTR或归因收入换一个名字。

观察数据只有在处理前混杂得到充分控制、处理定义一致、相应状态具有支持且干扰得到处理等条件下，才能识别这种差异。Gordon等对其研究场景中的大规模广告实验与观察方法进行比较，发现丰富特征并不保证观察估计复现随机实验\scite{rs-a17}这不是说任何观察法永远无效，而是不能用更强预测器代替未满足的识别条件。该正式研究与早期白皮书的实验集合不同，本章不混用它们的数量或效果。

用户留出通常按用户随机分配政策并比较完整结果窗口，避免按实际点击、实际曝光或完成购买后再分组。地域实验可在城市等区域调整投放强度，用处理前数据检查可比性并提高精度；跨地域购买、媒介覆盖和预算重分配仍会破坏隔离。两者都应记录自然内容和其他广告的替代，否则只增加某广告账目而减少自然成交，可能被误报为总体增长。

\subsubsection{Ghost Ads投放机会对照实验}
\label{sec:rs-advertising-ghost}

全体用户留出中，大量用户原本没有机会获得广告，会稀释可检测信号；只比较实际曝光者与未曝光者，又会引入投放选择。\term{Ghost Ads}由平台为对照用户记录“按相同投放规则本会获得本广告”的潜在机会，处理用户在对应机会实际投放，从而支持机会对应的随机比较\scite{rs-a16}它是一种实验设计和日志能力，不是预测转化的网络模型。

实验先固定用户级随机分组、广告资格、竞价与机会判定规则以及结果窗口。处理组按规定参与并展示；对照组不接受本广告，但保留影子竞价或等价机会记录，识别本会展示的机会，同时说明实际替代内容。分组必须不影响用于选择机会的前置信息。处理后的点击、访问或频次不能用于重新挑选“可比用户”，否则随机化不再保护该子集。

尤其当系统根据历史响应自适应投放时，首个实验机会之前两组尚未受到本实验广告影响，可以在适当实现下保持对称；之后处理可能改变访问次数、点击和后续资格。原文补充材料专门讨论这种机会与频次边界。因此，不宜把两组各自“曝光至少三次”的用户直接拿来比较，也不能仅凭有ghost日志就认为所有后续次数条件化都有效。

\begin{example}[全体留出与机会对应的分母]
教学假设两组各随机分配10000名用户，各有2000名在未受处理影响的共同规则下进入潜在机会集合。处理组机会用户有130人转化，对照组机会用户有100人转化；两组其余8000人均有400人转化。机会集合的转化率差为$130/2000-100/2000=0.015$，即1.5个百分点；全体分组的差为$530/10000-500/10000=0.003$，即0.3个百分点。

两者分母不同而不矛盾：在该假设中，机会覆盖比例0.2乘条件差0.015等于全体差0.003。机会处理组相对于相同规模对照多30次转化；若相对对照多支出600元，增量CPA点估计为20元。实际分析还需区间估计、用户级聚类和样本量设计，不能把这组教学计数视为效果已被高精度证实。
\end{example}

参数来自实验设计而非模型训练，平台需要实现机会判定、稳定随机分桶、shadow记录及完整结果收集。在线机会标记在已有拍卖之上通常只增加有限计算，但影子参与、日志存储与跨设备关联有成本；离线统计至少随用户和机会记录数线性增长。必须检验控制组的ghost判定是否复现处理组真实可展示条件，包括预算、底价、资格和素材版本。

预算重分配尤其关键。不给对照组展示后，若把节省资金全部投给处理组，处理组机会和频次就可能增加，所估计的是这一重新分配政策的效果，而非原投放规模下单次展示效果。影子报价还不能随意改变其他广告的支付，否则对照市场也受到处理。机会对应提高比较针对性，但不会自动修复拍卖干扰、身份缺失或结果窗口不一致。

\Needspace{90mm}
\subsubsection{CausCF：用户与对象条件下的增量估计}
\label{sec:rs-advertising-causcf}

平均实验效果不足以决定具体用户和对象的选择。\term{CausCF}将协同过滤扩展到用户、对象和处理三个维度，通过共享因子预测不同处理条件下的购买，再比较两个预测\scite{rs-a29}原文的处理是推荐曝光；迁移到广告时，应重新定义“投放”及对照内容，并核验训练数据是否支持这项处理，而不是直接把推荐实验结论移植为广告增量。

样本$(u,i,a,y)$包含用户、对象、实际处理$a\in\{0,1\}$及结果窗口内是否购买。用户和对象的处理前属性经嵌入与MLP形成$\bm{p}_u,\bm{q}_i\in\mathbb R^d$，处理标识形成$\bm{d}_a$。预测评分与购买概率为
\begin{equation}
  f_{u,i,a}
    =\langle\bm{p}_u,\bm{q}_i\rangle
      +\langle\bm{p}_u,\bm{d}_a\rangle
      +\langle\bm{q}_i,\bm{d}_a\rangle,\qquad
  \widehat m_{u,i}(a)=\sigma(f_{u,i,a}).
  \label{eq:rs-advertising-causcf}
\end{equation}
用户--对象项描述原有偏好，另外两项共享用户及对象对处理的响应。事实样本上的二元交叉熵与参数$\ell_2$正则共同训练表示；未观察的另一个处理结果不补零。点击和停留若在处理之后发生，不能作为要调整的处理前混杂直接加入。

使用时固定同一$(u,i)$，分别输入$a=1$和$a=0$。本章采用概率差
$\widehat\Delta_{u,i}=\widehat m_{u,i}(1)-\widehat m_{u,i}(0)$，
该差值以概率为尺度，乘100后可按百分点报告。原文也以未过sigmoid的评分差呈现效应，二者须区分：评分差是logit尺度差，不等于概率增量。该结构在评分尺度的用户--对象项会相消，处理效应由用户--处理和对象--处理两项组成，不能表示任意复杂的三方交互。

教学假设对象甲在不投放与投放下的购买概率为0.30、0.34，对象乙为0.05、0.15。按投放后响应，甲优于乙；按概率增量，乙的0.10大于甲的0.04。若每次购买贡献毛利100元，两次机会的增量毛利估计为4元和10元，再与各自预期支付比较。这里仍是教学预测；没有识别依据时，只能称为条件预测差，不能据此宣称真实增量。

低秩共享使稀疏用户和对象能够借用相近样本，但因子分解本身不提供随机化。原文保留无干扰、条件可忽略性与正支持假设，并以用户浏览停止位置附近的断点比较检验局部效果。这要求阈值附近潜在结果连续、没有精确操纵和伴随的其他处理变化；用户主动停止浏览还可能携带偏好信息。局部断点证据不能替代所有用户和位置的实验验证。

若用户和对象表示已缓存，每个处理评分为$O(d)$，新增用户还需属性编码；纯ID形式的存储约为$O((|U|+|I|+2)d)$。两种预测共享参数，区间估计应保留其相关性，不能把两套独立误差条直接相减。增量排序的决策接口见第~\ref{chap:rs-ranking}~章，识别与验证见第~\ref{chap:rs-evaluation}~章；本模型补充的是条件效应表达，预算和价格仍由投放控制处理。

\Needspace{90mm}
\subsubsection{MOTTO：多种干预与多个结果的效应估计}
\label{sec:rs-advertising-motto}

广告体验可能有多种可选干预，而且同时影响购买、互动和流失。为每个处理--结果组合独立训练，容易浪费相关任务之间的证据；将全部任务完全共享，又可能压平真正不同的效应。\term{MOTTO}设置全局共享、跨处理共享、跨结果共享及组合专属专家，经门控路由形成每个处理和结果下的预测\scite{rs-a30}本节机制依据作者公开实现及KDD 2025机构出版记录，正式论文全文未取得，因此不补未经核验的理论结论。

一个样本$(\bm{x}_n,a_n,\bm{y}_n)$包含处理前特征、实际接受的处理$a_n\in\{0,\ldots,K-1\}$及$R$维结果。模型输出$R\times K$矩阵$\widehat{\bm M}(\bm{x}_n)$。全局专家对全部组合可见；固定某个结果、跨处理共用的专家学习该结果中可共享的部分；固定某个处理、跨结果共用的专家学习该处理相关的共同部分；专属专家保留单个组合的信息。每个输出门控只混合与自身相关的专家，再由对应预测头输出结果。

“处理共享”在这里是跨不同处理共享同一结果的表示，不是只属于某一个处理的专属塔。作者实现对每条样本用实际处理索引读取$\widehat{\bm M}_{:,a_n}$，与事实结果$\bm y_n$比较；连续结果的示例使用平方误差，二元结果路径使用相应概率输出和交叉熵。缺少其他处理的结果，不意味着其他输出应为零。

带选择性分布对齐的版本额外取出跨处理共享专家的表示，按样本的实际处理分组，使各组表示分布接近合并后的参考分布。公开合成示例使用能量距离，训练器另有Sinkhorn距离实现；对齐项与事实预测损失加权后用梯度更新。选择性意味着只对指定共享表示施加这种压力，不把所有处理专属信息一并抹去。对齐强度仍需验证，表征分布接近也不证明未观测混杂已经消失。

推断时同一特征输入得到全部条件预测，再按每个结果比较处理与基准：
\[
  \widehat\Delta_{r,a}(\bm{x})
    =\widehat M_{r,a}(\bm{x})-\widehat M_{r,0}(\bm{x}),
      \qquad a=1,\ldots,K-1.
\]
这给出效应向量，而非一个天然可加的总分。购买概率、互动次数和流失概率有不同单位、方向及业务代价，交给第~\ref{chap:rs-constraints}~章的权衡接口之前须明确目标和限制。

\begin{table}[htbp]
  \centering
  \begin{tabularx}{\linewidth}{@{}Xrrr@{}}
    \toprule
    教学处理 & 购买概率 & 互动次数 & 流失概率\\
    \midrule
    基准：无新增广告干预 & 0.10 & 2.0 & 0.05\\
    处理一：增加广告展示 & 0.14 & 2.2 & 0.09\\
    处理二：保持负载、调整创意 & 0.13 & 2.4 & 0.04\\
    \bottomrule
  \end{tabularx}
  \caption{MOTTO输出接口的教学预测，不是论文结果。购买和流失采用固定用户窗口，互动为该窗口内期望次数；不同窗口的数值不能直接比较。}
  \label{tab:rs-advertising-motto}
\end{table}

按表~\ref{tab:rs-advertising-motto}中的购买、互动、流失顺序，两项处理的预测效应分别为
\[
  (0.04,0.2,0.04),\qquad (0.03,0.4,-0.01).
\]
只看购买，处理一稍高；若要求流失不能增加，处理一不合格。最终决策还需支付、预算以及效应不确定性，而不是让门控自行决定业务目标。

若每个共享类别和每个组合均设置固定数量专家，专家数仍包含$KR$个专属组，输出至少有$KR$项；共享减少重复学习，不意味着计算与处理结果数量无关。全前向还承担所有专家和门控的代价，批内分布距离可能需要成对样本计算。处理数很多而部分处理样本稀少时，支持不足、方差和对齐误差会更加明显。

MMoE回答的是如何共享预测多个结果，MOTTO另有处理维度，并以同一结果的条件差作为使用接口。这个结构区别不能代替因果识别：多处理的条件可交换性、支持、处理定义一致性及干扰仍须检验。公开代码的机制阅读也不是运行复现，正式实验数字和理论保证不由本节外推。

\subsection{受限追踪下的广告效果测量}
\label{sec:rs-advertising-limited-tracking}

无法完整关联用户路径时，不能假装仍拥有逐人逐触点的全部数据。能够在允许范围内关联的触点与转化，可以进行受贡献限制的归因统计；只有渠道、时间和地域汇总时，可以拟合聚合响应；独立增量实验则可补充聚合模型的识别信息。选择哪种数据形态，都要先说明目标是信用核算还是投放增量。

\subsubsection{Differentially Private Ad Conversion Measurement：贡献限制下的归因统计}
\label{sec:rs-advertising-dp-measurement}

\term{Differentially Private Ad Conversion Measurement}研究归因规则、隐私相邻关系、贡献限制范围和执行位置如何组合，才能控制最终归因数据的敏感度\scite{rs-a31}它不训练响应预测器，而是规定怎样从允许处理的触点与转化生成受保护报表。差分隐私的一般定义见\BookChapterRef{chap:trust-privacy}{基础、理论和可信性卷的“隐私保护”章}。

相邻关系规定一次需要保护的变化是什么：一条转化、一次展示、某用户在一个广告主上的活动，还是整名用户的全部活动。贡献限制范围规定哪些事件共用额度。保护一次转化，不等于保护转化很多次的用户；把“用户--发布方”写成“用户”也会改变承诺。规则还必须说明限制发生在归因前的输入事件，还是归因后的加权触点--转化对。

归因前限制会影响哪些事件能够参与路径核算。归因后限制则先算信用，再删除或缩减超额信用，通常不因后续裁剪而自动回退到另一个触点。两者会产生不同账目，不能只看上限数字相同就认为统计等价。原文在限定的归因规则和相邻关系集合上给出配置分类；它并未证明任意全局学习的归因模型搭配任意裁剪都有效。

一个常见误区是“每个发布方的最后触点最多记一份信用，所以保护用户在一个发布方的全部活动也只影响一份”。新增该发布方的一组末触，可能从很多其他发布方各夺走一份信用；即使新增一侧最后只保留少量，其他发布方的损失仍可能随数量增长。原文用此类构造说明，后归因限制与用户--发布方相邻关系并非任意可配。敏感度必须针对完整归因、裁剪、汇总管道计算，不能只检查最后的求和。

\begin{example}[归因前后裁剪的区别]
教学假设同一用户在同一广告主的有效窗口内先见展示A、再点击B，随后发生两次各计1的转化。未限制时，首触渠道账目为$(2,0)$，末触为$(0,2)$。现取用户级贡献上限2。

若按原文所讨论的归因前事件限制方式，展示A和点击B各消耗一份输入额度，后续转化无法再参与，报表为$(0,0)$。若先完成归因、再限制该用户最多贡献2份转化信用，两次转化都可保留，末触账目仍为$(0,2)$。同一个数字2在前者限制输入事件，在后者限制信用总量，信息损失不同。

固定后一种用户级增删相邻关系，假设归因逐用户独立、公开渠道桶预先固定，每用户信用向量的$\ell_1$范数不超过2，则汇总向量的$\ell_1$敏感度至多2。以$\epsilon=1$发布时，可对各桶加入尺度为2的独立Laplace噪声。假设本次噪声恰为$(0.3,-0.8)$，则$(0,2)$发布为$(0.3,1.2)$，不再严格守恒。该数值是一个假设噪声实现，不是噪声被限制在这个范围。
\end{example}

若相邻关系改为用另一用户记录替换原用户，以上简单上界变为4；若金额也要发布，还要限制单次金额或每用户金额总量。负噪声计数可以按预设规则截为零，这种后处理不增加隐私消耗，却会改变偏差。重复发布日报、周报及不同窗口需要计算累计隐私开销，不能每张表都独立宣称拥有完整的新额度。

按用户和时间分组后，首末触核算及贡献计数可在线性扫描中完成；排序、跨源关联及按多个维度发布会增加成本。用户级限制也要求能在可信处理边界内协调同一用户的贡献，不能一面无法关联贡献，一面声称已实施全用户上限。输出仍是带裁剪与噪声的归因信用，不是对照实验增量。预算配置应考虑低计数桶的不确定性和系统性裁剪损失，而不是把负噪声后的偶然高ROAS直接当作扩量信号。

\subsubsection{Bayesian MMM：聚合观测下的媒体效果估计}
\label{sec:rs-advertising-mmm}

只有每周渠道花费和销量时，无法重建个体触点路径，但可以研究投放与聚合结果的时间关系。\term{贝叶斯媒体组合模型}（Bayesian media mix model，\term{Bayesian MMM}）用广告滞后和饱和函数变换媒体变量，再接入回归和先验，得到参数及回报的后验分布\scite{rs-a32}本节依据Google 2017技术论文，不附会未核验的会议归属。

样本是按周对齐的$(y_t,x_{t,1:J},\bm z_t)$，其中$y_t$为销售收入，$x_{t,j}$为渠道花费，$\bm z_t$为事先选定的季节、价格、促销等控制因素。需要记录建模窗口之前的投放，避免把开始几周的历史影响误归于当前花费。若价格或促销本身受广告影响，是否将其作为控制变量取决于要估计总效应还是直接效应，不能不加区分地控制所有相关变量。

原文采用归一化\term{广告存量变换}：
\begin{equation}
  a_{t,j}
   =\frac{\sum_{\ell=0}^{L-1}w_j(\ell)x_{t-\ell,j}}
           {\sum_{\ell=0}^{L-1}w_j(\ell)},\qquad
  H_j(a)=\frac{a^{\nu_j}}{K_j^{\nu_j}+a^{\nu_j}}.
  \label{eq:rs-advertising-mmm-transform}
\end{equation}
$L$是最大滞后长度。几何权重$w_j(\ell)=\alpha_j^\ell$、$0<\alpha_j<1$令影响当期最高；延迟权重$w_j(\ell)=\alpha_j^{(\ell-\theta_j)^2}$允许峰值出现在之后。$K_j>0$是半饱和点，$\nu_j>0$控制曲线形状。$\nu_j>1$时低花费区可以先凸后凹，不能把所有Hill曲线都称为处处边际递减。

先滞后再饱和的加性模型为
\begin{equation}
  y_t=\beta_0+\sum_{j=1}^{J}\beta_jH_j(a_{t,j})
           +\langle\bm\gamma,\bm z_t\rangle+\varepsilon_t,
   \qquad \varepsilon_t\sim\mathcal N(0,\sigma^2).
  \label{eq:rs-advertising-mmm-regression}
\end{equation}
这个写法没有渠道协同项，误差独立性和恒定方差也需要诊断。先饱和再滞后会得到另一模型，不能在计算中互换顺序。常用先验限制$\alpha_j$在$[0,1)$、$K_j,\nu_j$为正，并在确有依据时限制$\beta_j$非负；非负先验是建模判断，不是数据已经证明广告不会伤害销售。

训练将观测似然与先验相乘，通过MCMC等方法抽取联合参数样本。原文比较专用采样器与Stan实现；通用采样机制见\BookChapterRef{chap:pgm-sampling}{模型卷的“基于采样的近似推断”章}。验收包括链收敛、有效样本数、残差时间相关、后验预测与先验敏感性。渠道花费几乎同步时，多种参数组合都能解释销量，更多采样不会凭空解决可辨识性。

\begin{example}[两渠道的滞后与回报]
\label{ex:rs-advertising-mmm}
教学假设$L=3$、几何衰减$\alpha=0.5$，之前花费均为0；渠道甲四周花费为$(7,0,0,0)$，乙为$(0,7,0,0)$。除以权重和1.75后，存量分别为$(4,2,1,0)$和$(0,4,2,1)$。设$K=2,\nu=1$、$\beta_{\text{甲}}=12,\beta_{\text{乙}}=9$，收入和花费使用同一货币尺度，两个渠道的模型收入贡献分别为$(8,6,4,0)$和$(0,6,4.5,3)$。

若把甲第一周花费从7改为0，后两周贡献也消失，因此完整窗口的模型回报为$(8+6+4)/7=18/7\approx2.571$，而不是仅取第一周的$8/7$。乙完整贡献为13.5，回报为$13.5/7\approx1.929$。这些是给定模型的预测差，不是未经假设即可成立的实验增量。

甲的单期饱和响应为$12a/(2+a)$，对存量$a$的导数为$24/(2+a)^2$；$a=1$时为$8/3$，$a=4$时为$2/3$。这说明不同花费区域的边际响应不同，但它还不是对原始当周花费的边际回报，后者需经过全部滞后路径求导。
\end{example}

计算渠道ROAS时，只把指定变化窗口内该渠道的花费置零，其他花费和控制因素按比较设定保持不变，并累计至最后一期受滞后影响的结果。边际ROAS则对该窗口的花费施加小幅变化，例如按原时间形状增加1\%，用预测收入差除以新增支出。平均回报高不意味着下一元边际回报高，改变花费在时间上的分布也可能改变结果。

不确定性应从联合后验向下传递。延续例~\ref{ex:rs-advertising-mmm}，教学上用等概率三点后验
$(\beta_{\text{甲}},\beta_{\text{乙}})=(10,11),(12,9),(14,7)$，
其他参数固定。甲ROAS取$(15/7,18/7,21/7)$，乙取$(16.5/7,13.5/7,10.5/7)$；范围分别约为$[2.143,3]$与$[1.5,2.357]$。这是刻意简化的离散后验支持范围，不是用三个样本估出的可靠95\%区间。在当前相同投放形状和饱和位置下，三种参数情形有两种更支持给甲增加少量预算、一种更支持乙。

实际预算可最大化后验平均预测收入，也可报告每个后验样本下的最优配置及其分散程度；二者不等价。不能把各参数的边际中位数拼成一套“典型参数”后只算一次回报，因为这会丢失参数相关性和非线性。配置交给第~\ref{sec:rs-advertising-channels}~节执行前，还要检查可用预算档、渠道最小额和调整时延。

设有$T$期、$J$渠道、最大滞后$L$，一次直接构造响应及似然约为$O(TJL)$；几何递推和缓存可减少重复计算。总训练成本还乘以采样链数、迭代数及每次梯度或条件更新的代价，预算搜索则对每个候选配置重复后验预测。稀疏时间点、共同趋势、遗漏促销和渠道联动会限制因果解释；曲线在已见范围内拟合良好，也不保证大幅扩量后的预测成立。

\Needspace{80mm}
\subsubsection{实验信息下的媒体效果校准}
\label{sec:rs-advertising-mmm-calibration}

聚合数据对单个渠道的效果常常约束较弱，独立增量实验能够补充信息。\term{Media Mix Model Calibration With Bayesian Priors}通过重新参数化，把指定窗口的广告回报直接作为模型参数，再为它设置来自实验或领域知识的先验\scite{rs-a33}这比反复试调回归系数先验更直接，但实验与模型必须描述可比较的处理、花费和结果窗口。

以单地域加性模型说明。设$D_{\mathcal W,j}$是窗口$\mathcal W$中渠道$j$的总花费，$G_{\mathcal W,j}(\bm\phi_j)$是在固定滞后、饱和参数$\bm\phi_j$下，将该窗口花费置零造成的单位系数响应总差，并包含窗口之后的滞后。则
\begin{equation}
  R_{\mathcal W,j}
    =\frac{\beta_jG_{\mathcal W,j}(\bm\phi_j)}
           {D_{\mathcal W,j}},
  \qquad
  \beta_j
    =\frac{R_{\mathcal W,j}D_{\mathcal W,j}}
           {G_{\mathcal W,j}(\bm\phi_j)}.
  \label{eq:rs-advertising-mmm-reparameterization}
\end{equation}
在分母非零时，用$R_{\mathcal W,j}$替换原$\beta_j$作为未知量。给回报、滞后和饱和参数设置先验，再代入原观测似然联合推断。回归系数此时依赖形状参数，不能仍要求它与这些参数保持原先的独立先验。原文的地域层级模型还包含地域随机效应，相同原则下应保留其汇总和缩放。

校准前先登记实验估计对象。若实验只增加10\%花费，其回报是该局部变化的回报；若实验关闭渠道，其比较更接近模型的置零回报，两者在饱和曲线上不同。实验期间的价格、促销、地域、计费口径以及自然流量替代也应可比。短期实验的高回报可能来自较低花费区域，不能直接当成季度高花费水平的平均回报先验。

延续例~\ref{ex:rs-advertising-mmm}，假设甲的投放实验只观察第一、二周，第一周花费7元，模型在这两周的收入差为14元，短窗口回报为2；第三周还有4元滞后贡献，完整回报为$18/7$。若校准参数是这个短窗口回报，应让模型也只在这两周累计收入差，得到$G_{\mathrm{short}}=7/6$，故$\beta=6R_{\mathrm{short}}$。若把2直接当成完整窗口回报，就会错配目标。

进一步作教学假设，实验信息给短窗口回报一个等概率两点先验$\{1.8,2.2\}$，对应$\beta=\{10.8,13.2\}$。在给定季度聚合数据、其他参数均固定的简化比较中，若两者观测似然之比为$2:1$，后验概率变为$(2/3,1/3)$，后验平均$\beta$为11.6。对应同一个单次脉冲完整窗口的平均回报为$1.5\times11.6/7\approx2.486$。这只是可复算的贝叶斯更新；真正季度回报还要用季度实际花费及滞后重新计算，不能复制该脉冲数值。

原文允许用实验子窗口定义回报参数，也允许在实验能够代表完整建模窗口时校准完整窗口。若实验没有观察到滞后尾部，可以明确校准其已测短期参数，再依靠额外假设推断尾部；不能把未测到的尾部当成实验已证明。多个实验可按各自花费构造聚合回报先验，但需要联合不确定性：若实验共享对照、地域或观测，不能按相互独立处理后重复提高精度。

训练仍需检查后验预测和敏感性。若校准用的实验结果已经包含于观测似然，必须明确联合模型或数据拆分，避免同一证据被重复使用；强先验与聚合数据冲突时，应检查目标错配和模型错误，而不是只报告更窄区间。回报先验不会消除遗漏混杂、拍卖干扰、函数形式错误或未来扩量外推风险。

重新参数化本身只增加计算窗口响应差的成本，主要计算仍在后验采样和预算搜索。输出应同时保留先验来源、校准窗口、后验回报及相应预算敏感性。该文机构页标注2024年，本次取得的官方稿首页标注2024年1月；本章按该版本引用，不虚写会议。实验提供外部证据，聚合模型负责在假设下连接时间与花费水平，实际配置仍受渠道可执行接口限制。

\section{本章小结}
\label{sec:rs-advertising-summary}

广告在共同推荐架构上加入了资金、支付规则和竞争。候选先满足商品、素材、预算与频次资格，响应模型再沿明确的展示--点击--转化链路提供概率与价值。点击选择、标签延迟和字段失准影响不同环节，分别修正之后仍要核验成熟窗口和数据支持。真实价值、预测价值、出价、支付与实际支出承担不同职责，不能因量纲相近就互相替代。

分配与支付决定当前机会怎样成交，实际支付更新余额，余额又改变下一次资格。单次竞价、周期节奏和跨渠道预算控制由此相连，但各自使用不同状态和动作。对偶控制依赖响应和约束假设，轨迹学习依赖日志覆盖，生成规划还要面对未来状态不可实现和评价器外推；这些方法都不能省去执行端的硬预算及频控。市场参与者共同响应时，固定价格回放也不足以代表全面部署后的结果。

创意生成扩展表达候选，事实核验限定哪些内容可以投放，创意选择用有限测试流量积累版本证据。个性化不能改写商品事实，换素材不能清空商品和品牌的曝光历史，短期点击也不能取代转化质量与持续体验。混排、负载、营销感和拒绝交互继续沿用展示章的接口，广告特有的资格、价格与资金状态在本章补齐。

效果测量最终决定反馈能支持什么判断。首末触和路径模型分配已发生转化的信用，随机对照及有条件的观察识别回答投放增加了什么，条件效应模型再支持更细的机会选择。追踪受限时，贡献限制和噪声发布保护归因统计，媒体组合模型利用聚合响应，实验先验补充外部证据；任何一种输出都不能越过其隐私、支持、时间和干扰条件。

至此，本部分从用户与对象理解，经候选、评分、展示和约束，延伸到长期反馈、评价、搜索及广告。共同问题始终是：系统依据哪些信息作出什么行动，行动改变了哪些状态，反馈又能证明什么。只有让这三者在相同口径下衔接，推荐、搜索与广告的模型改进才有机会转化为可核验、可持续的系统改进。
